\chapter{Relative Mealy Evaluation and Polynomial Response Semantics}
\label{ch:relative-mealy-evaluation-polynomial-response}
\label{ch:relative-mealy-program-categories}
\label{ch:polynomial-relative-mealy-actions}

\section{Orientation: evaluating and currying execution}
\label{sec:relative-response-orientation}

The preceding chapter associated to every Mealy morphism
\[
    \Phi=(P,\delta_P,\omega_P):
    \mathbb A\rightsquigarrow\mathbb B
\]
a labelled execution category
\[
\begin{tikzcd}[column sep=large]
    \mathbb A^{\mathrm{op}}
    &
    \operatorname{Exec}(\Phi)
        \arrow[l,"I_\Phi"']
        \arrow[r,"O_\Phi"]
    &
    \mathbb B^{\mathrm{op}},
\end{tikzcd}
\]
together with its finite trajectory sheaf and Segal composition structure.
The present chapter adds a downstream semantic context and develops two
presentations of the resulting relative execution semantics.

Let
\[
    Y=(Y\xrightarrow{r}B_0,\sigma)
\]
be a right \(\mathbb B\)-action.  Evaluating the emitted
\(\mathbb B\)-arrow in \(Y\) couples a Mealy state with a downstream state and
produces the relative state object
\[
    S_{\Phi,Y}=P\times_{q_P,B_0,r}Y.
\]
The associated transition is
\[
    (a,x,y)
    \longmapsto
    \bigl(
        \delta_P(a,x),
        \sigma(\omega_P(a,x),y)
    \bigr).
\]
Its action category is the relative execution, or relative program, category
\[
    \mathsf{Prog}_{\Phi,Y}.
\]
This is the uncurried operational presentation.

Whenever the relevant dependent products exist, the same labelled one-step
data transpose to a total response profile
\[
    c_{\Phi,Y}:
    S_{\Phi,Y}
    \longrightarrow
    \mathcal M_{\mathbb A,\mathbb B,Y}(S_{\Phi,Y}),
\]
with generalized-element formula
\[
    c_{\Phi,Y}(x,y)(a)
    =
    \left(
        \omega_P(a,x),
        \left(
            \delta_P(a,x),
            \sigma(\omega_P(a,x),y)
        \right)
    \right).
\]
This is the curried polynomial presentation.

The chapter is organized around the equivalence
\[
\begin{tikzcd}[column sep=huge]
    \text{labelled relative action}
        \arrow[r,shift left=1.1ex,"\text{dependent-product transpose}"]
    &
    \text{multiplicative response coalgebra}
        \arrow[l,shift left=1.1ex,"\text{flattening}"'].
\end{tikzcd}
\]
The action category of the left-hand structure and the flattened execution
category of the right-hand structure are canonically the same internal
category.

There are therefore four mutually compatible descriptions:
\[
\begin{array}{c|c}
    \text{description} & \text{principal object} \\
    \hline
    \text{unevaluated execution} & \operatorname{Exec}(\Phi) \\
    \text{evaluated relative execution} & \mathsf{Prog}_{\Phi,Y} \\
    \text{curried total response} & c_{\Phi,Y} \\
    \text{finite relative trajectories} &
        \operatorname{Path}(\mathsf{Prog}_{\Phi,Y}).
\end{array}
\]
The logical interpretation of selected execution witnesses, total response
profiles, and finite trajectories is deferred to the following chapter.

The principal results are as follows.
\begin{enumerate}
    \item A right \(\mathbb B\)-action is a terminal-output Mealy morphism
    \[
        \widehat Y:\mathbb B\rightsquigarrow\mathbf 1.
    \]
    Its execution category is its action category.

    \item Relative execution is Mealy composition:
    \[
        \mathsf{Prog}_{\Phi,Y}
        \cong
        \operatorname{Exec}(\widehat Y\circ\Phi).
    \]
    Consequently, by the temporal pullback theorem,
    \[
        \mathsf{Prog}_{\Phi,Y}
        \cong
        \operatorname{Exec}(\Phi)
        \times_{\mathbb B^{\mathrm{op}}}
        \int_{\mathbb B}Y.
    \]

    \item Coalgebras for the relative response polynomial are exactly typed
    labelled one-step response maps.  Multiplicative coalgebras are exactly
    labelled relative actions.

    \item Flattening a multiplicative coalgebra recovers its action category,
    including the input and downstream labellings.

    \item A Mealy simulation does not generally induce a strict coalgebra map
    over the response interface.  It induces an output-lax profile
    comparison, whose flattening is the contravariant relative execution
    functor obtained by downstream evaluation.
\end{enumerate}

The ambient hypotheses are used locally.  Finite limits suffice for right
actions, action categories, downstream evaluation, relative execution, and
flattened one-step spans.  The polynomial transposes require the explicitly
stated dependent products; local cartesian closure is a sufficient global
hypothesis.  Regular images, coherent joins, and first-order Heyting
quantifiers are not used in this chapter except in remarks pointing forward
to the dynamic-logic chapter.

We retain the composition convention of the preceding chapters.  For
internally composable arrows
\[
    a:u\to v,
    \qquad
    b:v\to w,
\]
write
\[
    m_A(a,b)=b\circ a.
\]
Right actions process target-to-source: at a state over \(w\), the arrow
\(b\) acts first and \(a\) acts second.  The corresponding output convention
is
\[
    m_B(\alpha,\beta)=\beta\circ\alpha.
\]
All generalized-element calculations below denote the corresponding morphism
equalities on the indicated pullback objects.

\section{Right actions and action categories}
\label{sec:relative-right-actions-action-categories}

Throughout this section let
\[
\mathbb A=(A_0,A_1,s_A,t_A,e_A,m_A)
\]
be an internal category in \(\mathcal E\).

The use of actions as Eilenberg--Moore algebras over a slice, and the
corresponding subobject-fibration logic used later in the chapter, follows the
standard categorical-logic perspective on indexed predicates and their
quantifiers \cite{Jacobs1999CategoricalLogic}.

\subsection{Right internal actions}
\label{subsec:relative-right-actions}

\begin{definition}[Right internal action]
\label{def:relative-right-action}
A \textbf{right action} of \(\mathbb A\) consists of a morphism
\[
p:X\to A_0
\]
together with a map
\[
\rho:
A_1\times_{t_A,A_0,p}X\to X
\]
such that
\[
p\rho=s_A\pi_1.
\]
The typing square for the action domain is
\[
\begin{tikzcd}
A_1\times_{t_A,A_0,p}X
\arrow[r,"\pi_X"]
\arrow[d,"\pi_1"']
&
X
\arrow[d,"p"]
\\
A_1
\arrow[r,"t_A"']
&
A_0
\arrow[ul,phantom,very near start,"\lrcorner"]
\end{tikzcd}
.
\]
The action map fits into
\[
\begin{tikzcd}
A_1\times_{t_A,A_0,p}X
\arrow[rr,"\rho"]
\arrow[dr,"s_A\pi_1"']
&&
X
\arrow[dl,"p"]
\\
& A_0 &
\end{tikzcd}
\]
and satisfies the unit and multiplication laws
\[
\rho(e_A(px),x)=x
\]
and
\[
\rho(b\circ a,x)=\rho(a,\rho(b,x))
\]
for internally composable arrows
\[
a:u\to v,
\qquad
b:v\to w,
\qquad
p(x)=w.
\]
Equivalently, using the multiplication map,
\[
\rho(m_A(a,b),x)=\rho(a,\rho(b,x)).
\]
\end{definition}

\begin{definition}[Equivariant map]
\label{def:relative-equivariant-map}
Let
\[
(X\xrightarrow{p}A_0,\rho_X)
\qquad\text{and}\qquad
(Z\xrightarrow{q}A_0,\rho_Z)
\]
be right \(\mathbb A\)-actions. An \(\mathbb A\)-equivariant map
\[
f:X\to Z
\]
is a morphism in \(\mathcal E\) satisfying
\[
qf=p
\]
and
\[
f(\rho_X(a,x))=\rho_Z(a,f(x)).
\]
Equivalently, the diagrams
\[
\begin{tikzcd}
X \arrow[rr,"f"] \arrow[dr,"p"'] && Z \arrow[dl,"q"] \\
& A_0 &
\end{tikzcd}
\]
and
\[
\begin{tikzcd}
A_1\times_{t_A,A_0,p}X
\arrow[r,"1\times f"]
\arrow[d,"\rho_X"']
&
A_1\times_{t_A,A_0,q}Z
\arrow[d,"\rho_Z"]
\\
X
\arrow[r,"f"']
&
Z
\end{tikzcd}
\]
commute.
\end{definition}

\begin{definition}[Category of right actions]
\label{def:relative-action-category-external}
The \textbf{category of right \(\mathbb A\)-actions}, denoted
\[
\mathbf{Act}(\mathbb A),
\]
has right internal \(\mathbb A\)-actions as objects and equivariant maps as
morphisms.
\end{definition}

\subsection{Actions as Eilenberg--Moore algebras}
\label{subsec:relative-actions-em}

\begin{proposition}[Actions as Eilenberg--Moore algebras]
\label{prop:relative-actions-em}
The internal category \(\mathbb A\), regarded as a monad in
\(\mathbf{Span}(\mathcal E)\), induces a monad
\[
T_{\mathbb A}:\mathcal E/A_0\to\mathcal E/A_0
\]
defined by
\[
T_{\mathbb A}(X\xrightarrow{p}A_0)
=
\left(
A_1\times_{t_A,A_0,p}X
\xrightarrow{s_A\pi_1}
A_0
\right).
\]
The category \(\mathbf{Act}(\mathbb A)\) is the Eilenberg--Moore category of
\(T_{\mathbb A}\).
\end{proposition}

\begin{proof}
\leavevmode
\begin{enumerate}
\item \textbf{Identify the endofunctor.}
For an object \(p:X\to A_0\) of \(\mathcal E/A_0\), composing the span
\[
A_0\xleftarrow{s_A}A_1\xrightarrow{t_A}A_0
\]
with \(p\) gives the pullback
\[
\begin{tikzcd}
A_1\times_{t_A,A_0,p}X
\arrow[r,"\pi_X"]
\arrow[d,"\pi_1"']
&
X
\arrow[d,"p"]
\\
A_1
\arrow[r,"t_A"']
&
A_0
\arrow[ul,phantom,very near start,"\lrcorner"]
\end{tikzcd}
,
\]
equipped with the structure map \(s_A\pi_1\) to \(A_0\).

\item \textbf{Unpack an algebra.}
A \(T_{\mathbb A}\)-algebra structure on \(p:X\to A_0\) is a map
\[
A_1\times_{t_A,A_0,p}X\to X
\]
over \(A_0\). This is precisely a map \(\rho\) satisfying
\[
p\rho=s_A\pi_1.
\]

\item \textbf{Compare the laws.}
The unit of the monad is induced by \(e_A:A_0\to A_1\), and the multiplication
is induced by \(m_A:A_1\times_{A_0}A_1\to A_1\). The Eilenberg--Moore unit
and multiplication equations are exactly
\[
\rho(e_A(px),x)=x
\]
and
\[
\rho(b\circ a,x)=\rho(a,\rho(b,x)).
\]

\item \textbf{Compare the morphisms.}
A morphism of Eilenberg--Moore algebras is a map over \(A_0\) commuting with
the algebra structures. This is exactly an equivariant map.
\end{enumerate}
\end{proof}

\subsection{Action categories}
\label{subsec:relative-action-categories}

Let
\[
X=(X\xrightarrow{p}A_0,\rho)
\]
be a right \(\mathbb A\)-action.

\begin{definition}[Action category]
\label{def:relative-action-category}
The \textbf{action category} of \(X\), denoted
\[
\int_{\mathbb A}X,
\]
is the internal category with
\[
\left(\int_{\mathbb A}X\right)_0=X,
\qquad
\left(\int_{\mathbb A}X\right)_1
=
A_1\times_{t_A,A_0,p}X.
\]
The source and target maps are
\[
s(a,x)=x,
\qquad
t(a,x)=\rho(a,x).
\]
They form the span
\[
\begin{tikzcd}
& A_1\times_{t_A,A_0,p}X
\arrow[dl,"\pi_X"']
\arrow[dr,"\rho"]
&
\\
X && X
\end{tikzcd}
.
\]
The identity map is
\[
e(x)=(e_A(px),x).
\]
Composition is given by
\[
\bigl((b,x),(a,\rho(b,x))\bigr)
\longmapsto
(m_A(a,b),x).
\]
Equivalently, the composite of
\[
x\xrightarrow{\,b\,}\rho(b,x)
\xrightarrow{\,a\,}\rho(a,\rho(b,x))
\]
is labelled by
\[
b\circ a=m_A(a,b).
\]
\end{definition}

\begin{theorem}[Action categories]
\label{thm:relative-action-category-internal}
For every right \(\mathbb A\)-action \(X\), the data of
Definition~\ref{def:relative-action-category} define an internal category
\[
\int_{\mathbb A}X.
\]
Equivalently, the span
\[
X
\xleftarrow{\pi_X}
A_1\times_{t_A,A_0,p}X
\xrightarrow{\rho}
X
\]
is a monoid object in
\[
\operatorname{EndSpan}_{\mathcal E}(X).
\]
\end{theorem}

\begin{proof}
\leavevmode
\begin{enumerate}
\item \textbf{Check the target map.}
The equation \(p\rho=s_A\pi_1\) says that the target of an action arrow is
typed over the source of its input arrow. If
\[
a:u\to v,
\qquad
p(x)=v,
\]
then \(\rho(a,x)\) lies over \(u\).

\item \textbf{Check identities.}
The identity arrow at \(x\) is
\[
(e_A(px),x).
\]
It lies in the pullback because \(t_Ae_Ap=p\). The action unit law gives
\[
\rho(e_A(px),x)=x.
\]

\item \textbf{Describe composable pairs.}
A composable pair in \(\int_{\mathbb A}X\) has the form
\[
\bigl((b,x),(a,\rho(b,x))\bigr)
\]
with
\[
a:u\to v,
\qquad
b:v\to w,
\qquad
p(x)=w.
\]
The corresponding pullback diagram is
\[
\begin{tikzcd}
\left(A_1\times_{A_0}X\right)
\times_X
\left(A_1\times_{A_0}X\right)
\arrow[r]
\arrow[d]
&
A_1\times_{A_0}X
\arrow[d,"\pi_X"]
\\
A_1\times_{A_0}X
\arrow[r,"\rho"']
&
X
\arrow[ul,phantom,very near start,"\lrcorner"]
\end{tikzcd}
.
\]

\item \textbf{Check composition.}
The composite
\[
(m_A(a,b),x)
\]
is defined because
\[
t_A(m_A(a,b))=t_A(b)=p(x).
\]
Its target is
\[
\rho(m_A(a,b),x)=\rho(a,\rho(b,x))
\]
by the action multiplication law.

\item \textbf{Verify the laws.}
The identity and associativity laws follow from the unit and associativity
laws of \(\mathbb A\) together with the action laws.
\end{enumerate}
\end{proof}

\subsection{Input-discrete functors}
\label{subsec:relative-input-discrete-functors}

The action category has a canonical internal functor
\[
\pi_X:\int_{\mathbb A}X\to\mathbb A^{\mathrm{op}}
\]
whose object part is
\[
p:X\to A_0
\]
and whose arrow part is the projection
\[
\pi_1:A_1\times_{t_A,A_0,p}X\to A_1.
\]
The source square of this functor is the defining pullback
\[
\begin{tikzcd}
A_1\times_{t_A,A_0,p}X
\arrow[r,"\pi_1"]
\arrow[d,"\pi_X"']
&
A_1
\arrow[d,"t_A"]
\\
X
\arrow[r,"p"']
&
A_0
\arrow[ul,phantom,very near start,"\lrcorner"]
\end{tikzcd}
.
\]
The use of \(t_A\) rather than \(s_A\) is because the codomain is
\(\mathbb A^{\mathrm{op}}\).

\begin{definition}[Input-discrete functor]
\label{def:relative-input-discrete}
Let
\[
J:\mathbb X\to\mathbb A^{\mathrm{op}}
\]
be an internal functor with object map
\[
J_0:X_0\to A_0.
\]
We say that \(J\) is \textbf{input-discrete} if the square
\[
\begin{tikzcd}
X_1 \arrow[r,"J_1"] \arrow[d,"s_X"']
&
A_1 \arrow[d,"t_A"]
\\
X_0 \arrow[r,"J_0"']
&
A_0
\end{tikzcd}
\]
is a pullback.
\end{definition}

\begin{theorem}[Actions as input-discrete functors]
\label{thm:relative-actions-input-discrete}
To give a right \(\mathbb A\)-action
\[
(X\xrightarrow{p}A_0,\rho)
\]
is equivalently to give an internal category \(\mathbb X\) with object of
objects \(X\), together with an input-discrete internal functor
\[
J:\mathbb X\to\mathbb A^{\mathrm{op}}
\]
whose object map is \(p\).
\end{theorem}

\begin{proof}
\leavevmode
\begin{enumerate}
\item \textbf{Construct the functor from an action.}
Given a right action \(X\), take
\[
\mathbb X=\int_{\mathbb A}X.
\]
Its arrow object is
\[
A_1\times_{t_A,A_0,p}X,
\]
and the functor to \(\mathbb A^{\mathrm{op}}\) is given by the projections
\[
p:X\to A_0,
\qquad
\pi_1:A_1\times_{t_A,A_0,p}X\to A_1.
\]
The input-discrete square is the pullback defining the arrow object.

\item \textbf{Recover an action from an input-discrete functor.}
Conversely, suppose
\[
J:\mathbb X\to\mathbb A^{\mathrm{op}}
\]
is input-discrete with object map \(p:X\to A_0\). The pullback condition gives
a canonical isomorphism
\[
X_1\cong A_1\times_{t_A,A_0,p}X.
\]
Transporting the target map of \(\mathbb X\) across this isomorphism gives
\[
\rho:A_1\times_{t_A,A_0,p}X\to X.
\]

\item \textbf{Check the typing.}
The target compatibility of \(J\) gives
\[
p\rho=s_A\pi_1.
\]
This is the required action typing equation.

\item \textbf{Check the laws.}
Under the pullback identification, the identity map of \(\mathbb X\) is
\[
x\longmapsto(e_A(px),x)
\]
because \(J\) preserves identities. Hence the unit law for \(\rho\) follows
from the identity law in \(\mathbb X\). Composition preservation by \(J\)
forces the transported composition map to be
\[
\bigl((b,x),(a,\rho(b,x))\bigr)
\longmapsto
(m_A(a,b),x).
\]
Associativity in \(\mathbb X\) gives
\[
\rho(m_A(a,b),x)=\rho(a,\rho(b,x)).
\]

\item \textbf{Compare the two constructions.}
The two assignments are inverse up to the canonical pullback isomorphism.
\end{enumerate}
\end{proof}


\subsection{Right actions as terminal-output Mealy morphisms}
\label{subsec:relative-actions-terminal-output-mealy}

Let
\[
    Y=(Y\xrightarrow{r}B_0,\sigma)
\]
be a right \(\mathbb B\)-action.  Regard the terminal internal category
\(\mathbf 1\) as having one object and one arrow.

\begin{definition}[Terminal-output Mealy morphism of an action]
\label{def:relative-terminal-output-mealy-action}
The \textbf{terminal-output Mealy morphism} associated to \(Y\) is
\[
    \widehat Y:\mathbb B\rightsquigarrow\mathbf 1
\]
with carrier span
\[
\begin{tikzcd}
    & Y \arrow[dl,"r"'] \arrow[dr] & \\
    B_0 && 1,
\end{tikzcd}
\]
transition map
\[
    \delta_{\widehat Y}(\beta,y):=\sigma(\beta,y),
\]
and the unique output map to the arrow object of \(\mathbf 1\).
\end{definition}

\begin{proposition}[Actions are terminal-output Mealy morphisms]
\label{prop:relative-actions-terminal-output-mealy}
The Mealy unit and multiplication laws for \(\widehat Y\) are exactly the unit
and multiplication laws of the right \(\mathbb B\)-action \(Y\).
Conversely, every Mealy morphism \(\mathbb B\rightsquigarrow\mathbf 1\) is
canonically a right \(\mathbb B\)-action.
\end{proposition}

\begin{proof}
The carrier anchor to \(B_0\) supplies the object of the action.  Since the
output category is terminal, the output map and all of its axioms are unique.
The transition map has domain
\[
    B_1\times_{t_B,B_0,r}Y
\]
and lies over \(s_B\).  The Mealy identity equation is
\[
    \delta_{\widehat Y}(e_B(r(y)),y)=y,
\]
which is the action unit law.  For composable arrows
\[
    \alpha:b_0\to b_1,
    \qquad
    \beta:b_1\to b_2,
\]
and \(y\) over \(b_2\), the Mealy multiplication equation is
\[
    \delta_{\widehat Y}(\beta\circ\alpha,y)
    =
    \delta_{\widehat Y}
    \bigl(\alpha,\delta_{\widehat Y}(\beta,y)\bigr),
\]
which is the right-action multiplication law.  The converse reads the same
equations in reverse.
\end{proof}

\begin{proposition}[Execution category of an action]
\label{prop:relative-exec-terminal-output-action}
There is a canonical isomorphism of internal categories
\[
    \operatorname{Exec}(\widehat Y)
    \cong
    \int_{\mathbb B}Y.
\]
Under this isomorphism, the input-labelling functor
\[
    I_{\widehat Y}:\operatorname{Exec}(\widehat Y)
    \to\mathbb B^{\mathrm{op}}
\]
is the standard projection of the action category to
\(\mathbb B^{\mathrm{op}}\).
\end{proposition}

\begin{proof}
Both internal categories have object object \(Y\) and arrow object
\[
    B_1\times_{t_B,B_0,r}Y.
\]
In both cases the source is the current state \(y\), the target is
\(\sigma(\beta,y)\), the identity is \((e_B(r(y)),y)\), and composition is
induced by \(m_B\) together with the action multiplication law.  The
input-labelling arrow map is the projection to \(B_1\), viewed in
\(\mathbb B^{\mathrm{op}}\).
\end{proof}

\section{Downstream evaluation of output-comparison cells}
\label{sec:relative-kleisli-evaluation}

The construction in this section treats an output-comparison Kleisli cell as a
context-dependent comparison which becomes an actual map only after a
downstream action has been chosen.

Fix internal categories
\[
\mathbb A=(A_0,A_1,s_A,t_A,e_A,m_A),
\qquad
\mathbb B=(B_0,B_1,s_B,t_B,e_B,m_B).
\]
Let
\[
\mathcal C_{A,B}:=\mathbf{Span}(\mathcal E)(A_0,B_0)
\]
be the local span category. The output-comparison monad is
\[
T_B:=B\circ -:\mathcal C_{A,B}\to\mathcal C_{A,B}.
\]
For a span
\[
P=\left(A_0\xleftarrow{p_P}P\xrightarrow{q_P}B_0\right),
\]
the object \(T_B(P)=B\circ P\) has apex
\[
P\times_{q_P,B_0,s_B}B_1
\]
and fits into
\[
\begin{tikzcd}
P\times_{q_P,B_0,s_B}B_1
\arrow[r,"\pi_{B_1}"]
\arrow[d,"\pi_P"']
&
B_1
\arrow[d,"s_B"]
\\
P
\arrow[r,"q_P"']
&
B_0
\arrow[ul,phantom,very near start,"\lrcorner"]
\end{tikzcd}
.
\]

Let
\[
Y=(Y\xrightarrow{r}B_0,\sigma)
\]
be a right \(\mathbb B\)-action.

\subsection{Evaluation on carriers}
\label{subsec:relative-evaluation-carriers}

\begin{definition}[Downstream evaluation on carriers]
\label{def:relative-downstream-evaluation-carriers}
For a span
\[
P=
\left(
A_0\xleftarrow{p_P}P\xrightarrow{q_P}B_0
\right)
\]
define
\[
\operatorname{Ev}_Y(P)
:=
P\times_{q_P,B_0,r}Y
\xrightarrow{p_P\pi_P}
A_0.
\]
The underlying pullback is
\[
\begin{tikzcd}
\operatorname{Ev}_Y(P)
\arrow[r,"\pi_Y"]
\arrow[d,"\pi_P"']
&
Y
\arrow[d,"r"]
\\
P
\arrow[r,"q_P"']
&
B_0
\arrow[ul,phantom,very near start,"\lrcorner"]
\end{tikzcd}
.
\]
\end{definition}

\subsection{Evaluation of Kleisli comparisons}
\label{subsec:relative-evaluation-kleisli-comparisons}

\begin{definition}[Evaluation of a Kleisli comparison]
\label{def:relative-kleisli-comparison-evaluation}
Let
\[
\Gamma:Q\to T_B(P)=B\circ P
\]
be a Kleisli arrow in \(\mathrm{Kl}(T_B)\). In components, write
\[
\Gamma(u)=(\gamma u,\lambda_u),
\]
where
\[
\gamma:Q\to P,
\qquad
\lambda:Q\to B_1,
\]
and
\[
p_P\gamma=p_Q,
\qquad
s_B\lambda=q_P\gamma,
\qquad
t_B\lambda=q_Q.
\]
The component data are displayed by
\[
\begin{tikzcd}[column sep=large,row sep=large]
&& B_1
  \arrow[dd,bend left=18,"s_B"]
  \arrow[dd,bend right=18,"t_B"']
&
\\
Q
  \arrow[rru,"\lambda"]
  \arrow[r,"\gamma"']
  \arrow[drr,"q_Q"']
  \arrow[ddrr,bend right=16,"p_Q"']
&
P
  \arrow[dr,"q_P"]
  \arrow[rdd,bend left=16,"p_P"]
&&
\\
&& B_0 &
\\
&& A_0 &
\end{tikzcd}
.
\]
Define
\[
\operatorname{Ev}_Y(\Gamma):
Q\times_{q_Q,B_0,r}Y
\longrightarrow
P\times_{q_P,B_0,r}Y
\]
by
\[
\operatorname{Ev}_Y(\Gamma)(u,z)
=
\bigl(\gamma u,\sigma(\lambda_u,z)\bigr).
\]
\end{definition}

\begin{proposition}[Evaluation is functorial on output-comparison cells]
\label{prop:relative-evaluation-functorial}
For every right \(\mathbb B\)-action \(Y\), the construction
\[
P\longmapsto \operatorname{Ev}_Y(P),
\qquad
\Gamma\longmapsto \operatorname{Ev}_Y(\Gamma)
\]
defines a functor
\[
\operatorname{Ev}_Y:\mathrm{Kl}(T_B)\to \mathcal E/A_0.
\]
\end{proposition}

\begin{proof}
\leavevmode
\begin{enumerate}
\item \textbf{Check well-definedness.}
If
\[
r(z)=q_Q(u)=t_B(\lambda_u),
\]
then \(\sigma(\lambda_u,z)\) is defined. Moreover
\[
r\sigma(\lambda_u,z)=s_B(\lambda_u)=q_P(\gamma u),
\]
so
\[
(\gamma u,\sigma(\lambda_u,z))
\]
lies in \(P\times_{B_0}Y\). It lies over \(A_0\) because
\[
p_P\gamma=p_Q.
\]

\item \textbf{Check identities.}
The Kleisli identity on \(P\) is
\[
x\longmapsto (x,e_B(q_Px)).
\]
Evaluating gives
\[
(x,z)\longmapsto
(x,\sigma(e_B(q_Px),z))
=
(x,z)
\]
by the unit law of the action \(Y\).

\item \textbf{Compute Kleisli composition.}
Let
\[
R\xrightarrow{\Delta}T_B(Q),
\qquad
Q\xrightarrow{\Gamma}T_B(P)
\]
be Kleisli arrows, with
\[
\Delta(w)=(\eta w,\mu_w),
\qquad
\Gamma(q)=(\gamma q,\lambda_q).
\]
The Kleisli composite has component
\[
w\longmapsto
\bigl(\gamma\eta w,\;m_B(\lambda_{\eta w},\mu_w)\bigr).
\]
In arrow notation this is the composite
\[
q_P(\gamma\eta w)
\xrightarrow{\lambda_{\eta w}}
q_Q(\eta w)
\xrightarrow{\mu_w}
q_R(w).
\]

\item \textbf{Evaluate the composite.}
Evaluation at \((w,z)\) gives
\[
\bigl(\gamma\eta w,
\sigma(m_B(\lambda_{\eta w},\mu_w),z)\bigr).
\]
By the right-action multiplication law, this equals
\[
\bigl(\gamma\eta w,
\sigma(\lambda_{\eta w},\sigma(\mu_w,z))\bigr),
\]
which is
\[
\operatorname{Ev}_Y(\Gamma)
\bigl(
\operatorname{Ev}_Y(\Delta)(w,z)
\bigr).
\]
\end{enumerate}
\end{proof}


\section{Relative execution as Mealy composition}
\label{sec:relative-execution-as-composition}

Fix a Mealy morphism
\[
    \Phi=(P,\delta_P,\omega_P):
    \mathbb A\rightsquigarrow\mathbb B
\]
and a right \(\mathbb B\)-action \(Y\).  Form the composite
\[
    \widehat Y\circ\Phi:
    \mathbb A\rightsquigarrow\mathbf 1.
\]
Its carrier is the span composite of the carrier of \(\Phi\) with the carrier
of \(\widehat Y\), hence its state object is
\[
    S_{\Phi,Y}
    :=
    P\times_{q_P,B_0,r}Y.
\]

\begin{proposition}[Transition of the evaluated composite]
\label{prop:relative-evaluated-composite-transition}
The transition of \(\widehat Y\circ\Phi\) is
\[
    (a,x,y)
    \longmapsto
    \bigl(
        \delta_P(a,x),
        \sigma(\omega_P(a,x),y)
    \bigr).
\]
\end{proposition}

\begin{proof}
Composition of Mealy morphisms first processes the input arrow in \(\Phi\),
producing the next \(P\)-state \(\delta_P(a,x)\) and the emitted
\(\mathbb B\)-arrow \(\omega_P(a,x)\).  The second machine
\(\widehat Y\) processes that arrow by the action map \(\sigma\).  The carrier
pullback condition is preserved because
\[
    r\sigma(\omega_P(a,x),y)
    =s_B\omega_P(a,x)
    =q_P\delta_P(a,x).
\]
\end{proof}

\begin{definition}[Relative execution category]
\label{def:relative-execution-category-composite}
The \textbf{relative execution category} of \((\Phi,Y)\) is
\[
    \mathsf{Prog}_{\Phi,Y}
    :=
    \operatorname{Exec}(\widehat Y\circ\Phi).
\]
The notation \(\mathsf{Prog}_{\Phi,Y}\) is retained because this internal
category is the program category used by the dynamic-logic semantics of the
following chapter.
\end{definition}

\begin{theorem}[Relative execution as a labelled pullback]
\label{thm:relative-execution-labelled-pullback}
There is a canonical isomorphism
\[
    \mathsf{Prog}_{\Phi,Y}
    \cong
    \operatorname{Exec}(\Phi)
    \times_{\mathbb B^{\mathrm{op}}}
    \int_{\mathbb B}Y.
\]
The two projections are respectively the input execution comparison and the
downstream comparison functor.
\end{theorem}

\begin{proof}
By Definition~\ref{def:relative-execution-category-composite},
\[
    \mathsf{Prog}_{\Phi,Y}
    =
    \operatorname{Exec}(\widehat Y\circ\Phi).
\]
The execution-composition theorem of the preceding chapter gives
\[
    \operatorname{Exec}(\widehat Y\circ\Phi)
    \cong
    \operatorname{Exec}(\Phi)
    \times_{\mathbb B^{\mathrm{op}}}
    \operatorname{Exec}(\widehat Y).
\]
Apply Proposition~\ref{prop:relative-exec-terminal-output-action} to identify
\(\operatorname{Exec}(\widehat Y)\) with \(\int_{\mathbb B}Y\).  The
statements about the two projections follow from the construction of the
labelled pullback.
\end{proof}

\begin{corollary}[Relative finite trajectories]
\label{cor:relative-finite-trajectories-pullback}
For every finite length \(n\),
\[
    \operatorname{Path}(\mathsf{Prog}_{\Phi,Y})(n)
    \cong
    \mathscr T_\Phi(n)
    \times_{\operatorname{Path}(\mathbb B^{\mathrm{op}})(n)}
    \operatorname{Path}\!\left(\int_{\mathbb B}Y\right)(n).
\]
These degreewise isomorphisms assemble into an isomorphism of labelled
trajectory objects.
\end{corollary}

\begin{proof}
The finite-path construction preserves finite limits.  Apply it to the
pullback of Theorem~\ref{thm:relative-execution-labelled-pullback}.
\end{proof}

The remaining operational sections calculate this composite explicitly and
establish its functoriality under downstream maps and Mealy simulations.  They
are retained because the explicit formulas are required later when the same
data are transposed into polynomial response profiles.

\section{Restriction of downstream actions along Mealy morphisms}
\label{sec:relative-evaluation-mealy-morphisms}

Let
\[
\Phi=(P,\delta_P,\omega_P):\mathbb A\rightsquigarrow\mathbb B
\]
be a Mealy morphism. Thus \(P\) is a span
\[
\begin{tikzcd}
& P \arrow[dl,"p_P"'] \arrow[dr,"q_P"] & \\
A_0 && B_0
\end{tikzcd}
,
\]
and
\[
\delta_P:A_1\times_{t_A,A_0,p_P}P\to P,
\qquad
\omega_P:A_1\times_{t_A,A_0,p_P}P\to B_1
\]
satisfy
\[
p_P\delta_P=s_A\pi_1,
\qquad
s_B\omega_P=q_P\delta_P,
\qquad
t_B\omega_P=q_P\pi_P,
\]
together with the Mealy unit and multiplication laws.

Let
\[
Y=(Y\xrightarrow{r}B_0,\sigma)
\]
be a right \(\mathbb B\)-action.

\subsection{The restricted action}
\label{subsec:relative-restricted-action}

\begin{definition}[Restricted downstream action]
\label{def:relative-restricted-action}
The \textbf{restriction of \(Y\) along \(\Phi\)} is the right
\(\mathbb A\)-action
\[
\Phi^\ast Y
\]
whose underlying object over \(A_0\) is
\[
S_{\Phi,Y}:=
P\times_{q_P,B_0,r}Y
\xrightarrow{p_P\pi_P}
A_0,
\]
and whose action map is
\[
A_1\times_{t_A,A_0,p_P\pi_P}S_{\Phi,Y}
\to S_{\Phi,Y}
\]
given by
\[
a\cdot(x,y)
=
\bigl(
\delta_P(a,x),
\sigma(\omega_P(a,x),y)
\bigr).
\]
The action square has the form
\[
\begin{tikzcd}
A_1\times_{A_0}S_{\Phi,Y}
\arrow[rr,"{(a,x,y)\mapsto(\delta_P(a,x),\sigma(\omega_P(a,x),y))}"]
\arrow[dr,"s_A\pi_1"']
&&
S_{\Phi,Y}
\arrow[dl,"p_P\pi_P"]
\\
& A_0 &
\end{tikzcd}
.
\]
\end{definition}

\begin{theorem}[Mealy restriction of downstream actions]
\label{thm:relative-mealy-restriction-actions}
For every Mealy morphism
\[
\Phi:\mathbb A\rightsquigarrow\mathbb B
\]
and every right \(\mathbb B\)-action \(Y\), the object
\(\Phi^\ast Y\) of Definition~\ref{def:relative-restricted-action} is a right
\(\mathbb A\)-action.
\end{theorem}

\begin{proof}
\leavevmode
\begin{enumerate}
\item \textbf{Check that the formula lands in the pullback.}
If
\[
t_B\omega_P(a,x)=q_P(x)=r(y),
\]
then \(\sigma(\omega_P(a,x),y)\) is defined. Moreover
\[
r\sigma(\omega_P(a,x),y)
=
s_B\omega_P(a,x)
=
q_P\delta_P(a,x),
\]
so the displayed pair lies in
\[
P\times_{B_0}Y.
\]

\item \textbf{Check the projection to \(A_0\).}
The projection of the result to \(A_0\) is
\[
p_P\delta_P(a,x)=s_A(a).
\]
Thus the action has the required target-to-source typing.

\item \textbf{Check the unit law.}
Using the Mealy unit laws and the unit law of \(Y\),
\[
e_A(p_Px)\cdot(x,y)
=
\bigl(
\delta_P(e_A(p_Px),x),
\sigma(\omega_P(e_A(p_Px),x),y)
\bigr)
=
(x,y).
\]

\item \textbf{Check the multiplication law.}
Let
\[
a:u\to v,
\qquad
b:v\to w,
\qquad
p_P(x)=w.
\]
Then
\[
(b\circ a)\cdot(x,y)
=
\bigl(
\delta_P(b\circ a,x),
\sigma(\omega_P(b\circ a,x),y)
\bigr).
\]
The Mealy multiplication laws give
\[
\delta_P(b\circ a,x)
=
\delta_P(a,\delta_P(b,x))
\]
and
\[
\omega_P(b\circ a,x)
=
\omega_P(b,x)\circ\omega_P(a,\delta_P(b,x)).
\]
Equivalently,
\[
\omega_P(b\circ a,x)
=
m_B\bigl(
\omega_P(a,\delta_P(b,x)),
\omega_P(b,x)
\bigr).
\]
The right-action multiplication law for \(Y\) gives
\[
(b\circ a)\cdot(x,y)
=
a\cdot\bigl(b\cdot(x,y)\bigr).
\]
\end{enumerate}
\end{proof}

\subsection{Functoriality in the downstream action}
\label{subsec:relative-functoriality-downstream-action}

\begin{proposition}[Functoriality in the downstream action]
\label{prop:relative-functoriality-downstream-action}
For fixed \(\Phi:\mathbb A\rightsquigarrow\mathbb B\), the assignment
\[
Y\longmapsto \Phi^\ast Y
\]
extends to a functor
\[
\Phi^\ast:\mathbf{Act}(\mathbb B)\to\mathbf{Act}(\mathbb A).
\]
If
\[
f:Y\to Z
\]
is a \(\mathbb B\)-equivariant map, then
\[
\Phi^\ast f:S_{\Phi,Y}\to S_{\Phi,Z}
\]
is given by
\[
(x,y)\longmapsto(x,f(y)).
\]
\end{proposition}

\begin{proof}
\leavevmode
\begin{enumerate}
\item \textbf{Check well-definedness.}
The square
\[
\begin{tikzcd}
S_{\Phi,Y}
\arrow[r,"1_P\times f"]
\arrow[d]
&
S_{\Phi,Z}
\arrow[d]
\\
P
\arrow[r,"1_P"']
&
P
\end{tikzcd}
\]
is induced by the pullback defining \(S_{\Phi,Y}\), because \(f\) lies over
\(B_0\).

\item \textbf{Check the map over \(A_0\).}
The \(P\)-component is unchanged, so the map lies over \(A_0\).

\item \textbf{Check equivariance.}
For an admissible input arrow \(a\),
\[
f(\sigma_Y(\omega_P(a,x),y))
=
\sigma_Z(\omega_P(a,x),f(y))
\]
by equivariance of \(f\) as a map of right \(\mathbb B\)-actions.

\item \textbf{Check identity and composition.}
Both follow from the corresponding equations for maps in \(\mathcal E\).
\end{enumerate}
\end{proof}

\section{Evaluation of Mealy simulations}
\label{sec:relative-evaluation-mealy-simulations}

Let
\[
\Phi=(P,\delta_P,\omega_P),
\qquad
\Psi=(Q,\delta_Q,\omega_Q)
:
\mathbb A\rightsquigarrow\mathbb B
\]
be Mealy morphisms.

\subsection{Component form}
\label{subsec:relative-simulation-component-form}

\begin{definition}[Mealy simulation, component form]
\label{def:relative-mealy-simulation-component}
A \textbf{Mealy simulation}
\[
\Theta:\Phi\Rightarrow\Psi
\]
is a Kleisli output-comparison cell
\[
Q\to B\circ P.
\]
In components, write
\[
\Theta(y)=(\theta y,\alpha_y),
\]
where
\[
\theta:Q\to P,
\qquad
\alpha:Q\to B_1,
\]
and
\[
p_P\theta=p_Q,
\qquad
s_B\alpha=q_P\theta,
\qquad
t_B\alpha=q_Q.
\]
The typing data are summarized by
\[
\begin{tikzcd}[column sep=large,row sep=large]
&& B_1
  \arrow[dd,bend left=18,"s_B"]
  \arrow[dd,bend right=18,"t_B"']
&
\\
Q
  \arrow[rru,"\alpha"]
  \arrow[r,"\theta"']
  \arrow[drr,"q_Q"']
  \arrow[ddrr,bend right=16,"p_Q"']
&
P
  \arrow[dr,"q_P"]
  \arrow[rdd,bend left=16,"p_P"]
&&
\\
&& B_0 &
\\
&& A_0 &
\end{tikzcd}
.
\]
The simulation equations are
\[
\theta(\delta_Q(a,y))
=
\delta_P(a,\theta y),
\]
and
\[
\alpha_y\circ \omega_P(a,\theta y)
=
\omega_Q(a,y)\circ \alpha_{\delta_Q(a,y)}.
\]
Equivalently,
\[
m_B\bigl(\omega_P(a,\theta y),\alpha_y\bigr)
=
m_B\bigl(\alpha_{\delta_Q(a,y)},\omega_Q(a,y)\bigr).
\]
\end{definition}

\subsection{Evaluated simulations}
\label{subsec:relative-evaluated-simulations}

\begin{definition}[Evaluated simulation map]
\label{def:relative-evaluated-simulation-map}
Let \(Y\) be a right \(\mathbb B\)-action. The \textbf{evaluation} of
\(\Theta\) at \(Y\) is the map
\[
h_{\Theta,Y}:S_{\Psi,Y}\to S_{\Phi,Y}
\]
given by
\[
h_{\Theta,Y}(y,z)
=
\bigl(\theta y,\sigma(\alpha_y,z)\bigr).
\]
It is induced by the diagram
\[
\begin{tikzcd}
S_{\Psi,Y}
\arrow[r,"h_{\Theta,Y}"]
\arrow[d]
&
S_{\Phi,Y}
\arrow[d]
\\
Q
\arrow[r,"\theta"']
&
P
\end{tikzcd}
\]
together with the action of \(\alpha_y\) on the \(Y\)-coordinate.
\end{definition}

\begin{theorem}[Simulations evaluate to equivariant maps]
\label{thm:relative-simulations-evaluate-equivariantly}
For every Mealy simulation
\[
\Theta:\Phi\Rightarrow\Psi
\]
and every right \(\mathbb B\)-action \(Y\), the map
\[
h_{\Theta,Y}:S_{\Psi,Y}\to S_{\Phi,Y}
\]
is an \(\mathbb A\)-equivariant map
\[
\Psi^\ast Y\to \Phi^\ast Y.
\]
\end{theorem}

\begin{proof}
\leavevmode
\begin{enumerate}
\item \textbf{Check well-definedness.}
Since
\[
r\sigma(\alpha_y,z)=s_B\alpha_y=q_P(\theta y),
\]
the pair
\[
(\theta y,\sigma(\alpha_y,z))
\]
lies in \(S_{\Phi,Y}\).

\item \textbf{Check the map over \(A_0\).}
The equation
\[
p_P\theta=p_Q
\]
gives
\[
p_P(\theta y)=p_Q(y).
\]

\item \textbf{Compute the left-hand equivariance composite.}
Let
\[
a:u\to v,
\qquad
p_Q(y)=v,
\qquad
r(z)=q_Q(y).
\]
Then
\[
h_{\Theta,Y}\bigl(a\cdot_\Psi(y,z)\bigr)
=
h_{\Theta,Y}
\bigl(
\delta_Q(a,y),
\sigma(\omega_Q(a,y),z)
\bigr),
\]
which is
\[
\bigl(
\theta\delta_Q(a,y),
\sigma(\alpha_{\delta_Q(a,y)},\sigma(\omega_Q(a,y),z))
\bigr).
\]
Using the action multiplication law for \(Y\), the second component is
\[
\sigma\bigl(
\omega_Q(a,y)\circ \alpha_{\delta_Q(a,y)},
z
\bigr).
\]

\item \textbf{Use the simulation law.}
The simulation law gives
\[
\omega_Q(a,y)\circ \alpha_{\delta_Q(a,y)}
=
\alpha_y\circ \omega_P(a,\theta y).
\]
Thus the second component becomes
\[
\sigma\bigl(
\alpha_y\circ \omega_P(a,\theta y),
z
\bigr).
\]
Using the action multiplication law again, this is
\[
\sigma(\omega_P(a,\theta y),\sigma(\alpha_y,z)).
\]

\item \textbf{Compare with the right-hand equivariance composite.}
The state equation
\[
\theta\delta_Q(a,y)=\delta_P(a,\theta y)
\]
gives
\[
h_{\Theta,Y}\bigl(a\cdot_\Psi(y,z)\bigr)
=
a\cdot_\Phi h_{\Theta,Y}(y,z).
\]
\end{enumerate}
\end{proof}

\begin{corollary}[Contravariance in simulations]
\label{cor:relative-contravariance-actions}
For fixed \(Y\in\mathbf{Act}(\mathbb B)\), the assignment
\[
\Phi\longmapsto\Phi^\ast Y
\]
is contravariant in Mealy simulations:
\[
\Theta:\Phi\Rightarrow\Psi
\quad\longmapsto\quad
h_{\Theta,Y}:\Psi^\ast Y\to\Phi^\ast Y.
\]
Thus it defines a functor
\[
\mathbf{Mly}(\mathcal E)(\mathbb A,\mathbb B)^{\mathrm{op}}
\longrightarrow
\mathbf{Act}(\mathbb A).
\]
Here \(\mathbf{Mly}(\mathcal E)(\mathbb A,\mathbb B)\) is the hom-category
with the displayed simulation direction of
Chapter~\ref{ch:spans-internal-categories-mealy}.
\end{corollary}

\begin{proof}
\leavevmode
\begin{enumerate}
\item \textbf{Check identities.}
The identity simulation uses the identity output arrows of \(\mathbb B\).
Evaluation sends these identity output arrows to identity transitions in \(Y\).

\item \textbf{Check composition.}
Composition of simulations is Kleisli composition. By
Proposition~\ref{prop:relative-evaluation-functorial}, downstream evaluation
preserves Kleisli composition.

\item \textbf{Check the direction.}
A simulation
\[
\Theta:\Phi\Rightarrow\Psi
\]
is represented by a Kleisli cell from the carrier of \(\Psi\) to the
output-comparison object of the carrier of \(\Phi\). Hence evaluation produces
a map
\[
\Psi^\ast Y\to\Phi^\ast Y.
\]
\end{enumerate}
\end{proof}

\section{Concrete form of the relative execution category}
\label{sec:relative-program-categories}

Fix a Mealy morphism
\[
\Phi=(P,\delta_P,\omega_P):\mathbb A\rightsquigarrow\mathbb B
\]
and a right \(\mathbb B\)-action
\[
Y=(Y\xrightarrow{r}B_0,\sigma).
\]

\subsection{State and one-step execution objects}
\label{subsec:relative-state-primitive-objects}

\begin{definition}[Relative state and one-step execution objects]
\label{def:relative-state-and-primitive}
The \textbf{relative state object} is
\[
S_{\Phi,Y}:=P\times_{q_P,B_0,r}Y.
\]
A generalized element is written
\[
(x,y),
\qquad
q_P(x)=r(y).
\]
The \textbf{relative one-step execution object} is
\[
M_{\Phi,Y}
:=
A_1\times_{t_A,A_0,p_P\pi_P}S_{\Phi,Y}.
\]
It is defined by the pullback
\[
\begin{tikzcd}
M_{\Phi,Y}
\arrow[r,"\pi_S"]
\arrow[d,"\pi_A"']
&
S_{\Phi,Y}
\arrow[d,"p_P\pi_P"]
\\
A_1
\arrow[r,"t_A"']
&
A_0
\arrow[ul,phantom,very near start,"\lrcorner"]
\end{tikzcd}
.
\]
A generalized element is written
\[
(a,x,y),
\qquad
t_A(a)=p_P(x),
\qquad
q_P(x)=r(y).
\]
\end{definition}

\subsection{The program category}
\label{subsec:relative-program-category-definition}

\begin{definition}[Concrete relative program category]
\label{def:relative-mealy-program-category}
The relative execution category introduced in
Definition~\ref{def:relative-execution-category-composite} has the
\textbf{action-category presentation}
\[
\mathsf{Prog}_{\Phi,Y}
:=
\int_{\mathbb A}\Phi^\ast Y.
\]
Explicitly,
\[
(\mathsf{Prog}_{\Phi,Y})_0=S_{\Phi,Y},
\qquad
(\mathsf{Prog}_{\Phi,Y})_1=M_{\Phi,Y}.
\]
The source and target maps are
\[
s_{\Phi,Y}(a,x,y)=(x,y),
\]
and
\[
t_{\Phi,Y}(a,x,y)
=
\bigl(
\delta_P(a,x),
\sigma(\omega_P(a,x),y)
\bigr).
\]
They form the one-step execution span
\[
\mathbf m_{\Phi,Y}
=
\left(
S_{\Phi,Y}
\xleftarrow{s_{\Phi,Y}}
M_{\Phi,Y}
\xrightarrow{t_{\Phi,Y}}
S_{\Phi,Y}
\right).
\]
The identity is
\[
e_{\Phi,Y}(x,y)=(e_A(p_Px),x,y).
\]
Composition is
\[
\bigl(
(b,x,y),
(a,\delta_P(b,x),\sigma(\omega_P(b,x),y))
\bigr)
\longmapsto
(m_A(a,b),x,y).
\]
\end{definition}

\begin{theorem}[Agreement of the relative execution presentations]
\label{thm:relative-program-categories}
For every Mealy morphism
\[
\Phi:\mathbb A\rightsquigarrow\mathbb B
\]
and every right \(\mathbb B\)-action \(Y\), the data of
Definition~\ref{def:relative-mealy-program-category} define an internal
category
\[
\mathsf{Prog}_{\Phi,Y}.
\]
This category is canonically isomorphic both to
\[
\operatorname{Exec}(\widehat Y\circ\Phi)
\]
and to
\[
\operatorname{Exec}(\Phi)
\times_{\mathbb B^{\mathrm{op}}}
\int_{\mathbb B}Y.
\]
Equivalently, the span
\[
\mathbf m_{\Phi,Y}
=
\left(
S_{\Phi,Y}
\xleftarrow{s_{\Phi,Y}}
M_{\Phi,Y}
\xrightarrow{t_{\Phi,Y}}
S_{\Phi,Y}
\right)
\]
is a monoid object in
\[
\operatorname{EndSpan}_{\mathcal E}(S_{\Phi,Y}).
\]
\end{theorem}

\begin{proof}
\leavevmode
\begin{enumerate}
\item \textbf{Use the restricted action.}
By Theorem~\ref{thm:relative-mealy-restriction-actions},
\(\Phi^\ast Y\) is a right \(\mathbb A\)-action.

\item \textbf{Apply the action-category construction.}
Theorem~\ref{thm:relative-action-category-internal} applied to this action
gives the internal category
\[
\int_{\mathbb A}\Phi^\ast Y.
\]

\item \textbf{Identify the displayed formulas.}
The source, target, identity, and composition maps in
Definition~\ref{def:relative-mealy-program-category} are precisely the
specialized action-category maps.

\item \textbf{Compare with the composite and pullback presentations.}
Proposition~\ref{prop:relative-evaluated-composite-transition} identifies
the action above with the transition of \(\widehat Y\circ\Phi\).
Theorem~\ref{thm:relative-execution-labelled-pullback} then gives the
labelled pullback presentation.
\end{enumerate}
\end{proof}

\subsection{Program functors induced by simulations}
\label{subsec:relative-program-functors-simulations}

\begin{definition}[Program functor induced by a simulation]
\label{def:relative-program-functor-simulation}
Let
\[
\Theta:\Phi\Rightarrow\Psi
\]
be a Mealy simulation with components
\[
\theta:Q\to P,
\qquad
\alpha:Q\to B_1.
\]
For every right \(\mathbb B\)-action \(Y\), define
\[
H_{\Theta,Y}:
\mathsf{Prog}_{\Psi,Y}
\to
\mathsf{Prog}_{\Phi,Y}
\]
on objects by
\[
H_{\Theta,Y}(y,z)
=
\bigl(\theta y,\sigma(\alpha_y,z)\bigr),
\]
and on arrows by
\[
H_{\Theta,Y}(a,y,z)
=
\bigl(a,\theta y,\sigma(\alpha_y,z)\bigr).
\]
The arrow part is induced by the pullback square
\[
\begin{tikzcd}
M_{\Psi,Y}
\arrow[r,"k_{\Theta,Y}"]
\arrow[d,"\pi_S"']
&
M_{\Phi,Y}
\arrow[d,"\pi_S"]
\\
S_{\Psi,Y}
\arrow[r,"h_{\Theta,Y}"']
&
S_{\Phi,Y}
\end{tikzcd}
\]
because both one-step execution objects are pulled back from \(A_1\).
\end{definition}

\begin{theorem}[Simulations induce contravariant program functors]
\label{thm:relative-simulations-program-functors}
The maps in Definition~\ref{def:relative-program-functor-simulation} define an
internal functor
\[
H_{\Theta,Y}:
\mathsf{Prog}_{\Psi,Y}
\to
\mathsf{Prog}_{\Phi,Y}.
\]
Consequently, for fixed \(Y\), the assignment
\[
\Phi\longmapsto\mathsf{Prog}_{\Phi,Y}
\]
extends to a functor
\[
\mathbf{Mly}(\mathcal E)(\mathbb A,\mathbb B)^{\mathrm{op}}
\longrightarrow
\mathbf{Cat}_{\mathrm{int}}(\mathcal E)/\mathbb A^{\mathrm{op}}.
\]
\end{theorem}

\begin{proof}
\leavevmode
\begin{enumerate}
\item \textbf{Identify the object map.}
The object map is
\[
h_{\Theta,Y}:S_{\Psi,Y}\to S_{\Phi,Y}.
\]

\item \textbf{Identify the arrow map.}
The arrow map is the induced map on the pullbacks defining primitive execution
objects:
\[
(a,y,z)\longmapsto
(a,\theta y,\sigma(\alpha_y,z)).
\]

\item \textbf{Check source and target preservation.}
Source preservation is immediate. Target preservation is exactly the
equivariance of \(h_{\Theta,Y}\), proved in
Theorem~\ref{thm:relative-simulations-evaluate-equivariantly}.

\item \textbf{Check identities and composition.}
These follow from functoriality of action categories with respect to
equivariant maps.

\item \textbf{Check the slice condition.}
The functor lands over \(\mathbb A^{\mathrm{op}}\) because the arrow map leaves
the \(A_1\)-component unchanged and the object map lies over \(A_0\).
\end{enumerate}
\end{proof}

\section{Input and downstream comparison functors}
\label{sec:relative-input-downstream-comparison}

Let
\[
\Phi=(P,\delta_P,\omega_P):\mathbb A\rightsquigarrow\mathbb B
\]
be a Mealy morphism, and let
\[
Y=(Y\xrightarrow{r}B_0,\sigma)
\]
be a right \(\mathbb B\)-action.

\subsection{Input projection}
\label{subsec:relative-input-projection}

\begin{definition}[Input projection]
\label{def:relative-input-projection}
The \textbf{input projection}
\[
I_{\Phi,Y}:\mathsf{Prog}_{\Phi,Y}\to\mathbb A^{\mathrm{op}}
\]
is defined on objects by
\[
I_{\Phi,Y}(x,y)=p_P(x),
\]
and on arrows by
\[
I_{\Phi,Y}(a,x,y)=a.
\]
It is displayed by
\[
\begin{tikzcd}
M_{\Phi,Y}
\arrow[r,"\pi_A"]
\arrow[d,"s_{\Phi,Y}"']
&
A_1
\arrow[d,"t_A"]
\\
S_{\Phi,Y}
\arrow[r,"p_P\pi_P"']
&
A_0
\end{tikzcd}
.
\]
\end{definition}

\subsection{Downstream comparison}
\label{subsec:relative-downstream-comparison}

\begin{definition}[Downstream comparison functor]
\label{def:relative-downstream-comparison}
The \textbf{downstream comparison functor}
\[
\Lambda_{\Phi,Y}:\mathsf{Prog}_{\Phi,Y}\to\int_{\mathbb B}Y
\]
is defined on objects by
\[
\Lambda_{\Phi,Y}(x,y)=y,
\]
and on arrows by
\[
\Lambda_{\Phi,Y}(a,x,y)
=
(\omega_P(a,x),y).
\]
The arrow map is typed by the pullback
\[
\begin{tikzcd}
M_{\Phi,Y}
\arrow[rr,"{(a,x,y)\mapsto(\omega_P(a,x),y)}"]
\arrow[dr,"\pi_Y\pi_S"']
&&
B_1\times_{t_B,B_0,r}Y
\arrow[dl,"\pi_Y"]
\\
& Y &
\end{tikzcd}
.
\]
\end{definition}

\begin{theorem}[Input and downstream functors]
\label{thm:relative-input-downstream-functors}
The input projection
\[
I_{\Phi,Y}:\mathsf{Prog}_{\Phi,Y}\to\mathbb A^{\mathrm{op}}
\]
is input-discrete. The downstream comparison
\[
\Lambda_{\Phi,Y}:\mathsf{Prog}_{\Phi,Y}\to\int_{\mathbb B}Y
\]
is an internal functor.
\end{theorem}

\begin{proof}
\leavevmode
\begin{enumerate}
\item \textbf{Check input-discreteness.}
The input projection is the canonical projection associated to the right
\(\mathbb A\)-action \(\Phi^\ast Y\). Hence it is input-discrete by
Theorem~\ref{thm:relative-actions-input-discrete}.

\item \textbf{Check that the input projection is a functor.}
The source square is the defining pullback. The target square commutes because
\[
p_P\delta_P(a,x)=s_A(a),
\]
which is the target equation for a functor into \(\mathbb A^{\mathrm{op}}\).
Identity and composition preservation follow from the identity and
multiplication laws of \(\mathbb A\).

\item \textbf{Check the downstream arrow map.}
The arrow map of \(\Lambda_{\Phi,Y}\) is well-defined because
\[
t_B\omega_P(a,x)=q_P(x)=r(y).
\]

\item \textbf{Check source and target preservation.}
In \(\int_{\mathbb B}Y\),
\[
s(\omega_P(a,x),y)=y,
\]
and
\[
t(\omega_P(a,x),y)=\sigma(\omega_P(a,x),y).
\]
These are the source and the \(Y\)-component of the target in
\(\mathsf{Prog}_{\Phi,Y}\).

\item \textbf{Check identities and composition.}
Identity preservation follows from the Mealy unit law for \(\omega_P\). For
composition, let
\[
a:u\to v,
\qquad
b:v\to w,
\qquad
p_P(x)=w.
\]
The composite of the two downstream arrows has \(B\)-label
\[
m_B\bigl(\omega_P(a,\delta_P(b,x)),\omega_P(b,x)\bigr)
=
\omega_P(b,x)\circ \omega_P(a,\delta_P(b,x)).
\]
By the Mealy multiplication law, this equals
\[
\omega_P(b\circ a,x).
\]
Thus \(\Lambda_{\Phi,Y}\) preserves composition.
\end{enumerate}
\end{proof}

\begin{corollary}[Combined labelling]
\label{cor:relative-combined-labelling}
Since \(\mathcal E\) has finite limits, the product internal category
\[
\mathbb A^{\mathrm{op}}\times\int_{\mathbb B}Y
\]
exists. Hence \(I_{\Phi,Y}\) and \(\Lambda_{\Phi,Y}\) assemble into an internal
functor
\[
L_{\Phi,Y}:
\mathsf{Prog}_{\Phi,Y}
\to
\mathbb A^{\mathrm{op}}\times\int_{\mathbb B}Y.
\]
On objects,
\[
L_{\Phi,Y}(x,y)=(p_Px,y),
\]
and on arrows,
\[
L_{\Phi,Y}(a,x,y)=\bigl(a,(\omega_P(a,x),y)\bigr).
\]
\end{corollary}

\begin{proof}
\leavevmode
\begin{enumerate}
\item \textbf{Use the product universal property.}
The functors \(I_{\Phi,Y}\) and \(\Lambda_{\Phi,Y}\) have common domain
\(\mathsf{Prog}_{\Phi,Y}\). The product internal category therefore supplies
a unique induced functor.

\item \textbf{Identify the components.}
The displayed object and arrow maps are the two projections of this induced
functor.
\end{enumerate}
\end{proof}

\section{Output-comparison transformations after evaluation}
\label{sec:relative-output-comparison-transformations}

Let
\[
\Theta:\Phi\Rightarrow\Psi
\]
be a Mealy simulation with components
\[
\theta:Q\to P,
\qquad
\alpha:Q\to B_1.
\]
Let \(Y\) be a right \(\mathbb B\)-action.

The induced functor
\[
H_{\Theta,Y}:
\mathsf{Prog}_{\Psi,Y}\to\mathsf{Prog}_{\Phi,Y}
\]
strictly preserves input labels:
\[
I_{\Phi,Y}H_{\Theta,Y}=I_{\Psi,Y}.
\]
The output-comparison component \(\alpha\) becomes a natural transformation in
the downstream action category.

\subsection{The evaluated downstream transformation}
\label{subsec:relative-downstream-transformation}

\begin{definition}[Evaluated downstream comparison transformation]
\label{def:relative-downstream-comparison-transformation}
Define
\[
\lambda_{\Theta,Y}:
\Lambda_{\Psi,Y}
\Rightarrow
\Lambda_{\Phi,Y}H_{\Theta,Y}
\]
by assigning to a state
\[
(y,z)\in S_{\Psi,Y}
\]
the arrow
\[
(\alpha_y,z)
:
z\to \sigma(\alpha_y,z)
\]
of the action category \(\int_{\mathbb B}Y\).
The component lies in
\[
B_1\times_{t_B,B_0,r}Y
\]
because
\[
t_B\alpha_y=q_Q(y)=r(z).
\]
Its source and target are
\[
\pi_Y(\alpha_y,z)=z,
\qquad
\sigma(\alpha_y,z).
\]
The relevant typing equations for the action-arrow object are
\[
r\pi_Y=t_B\pi_{B_1},
\qquad
r\sigma=s_B\pi_{B_1}.
\]
Equivalently, the component is represented by the diagram
\[
\begin{tikzcd}[column sep=large,row sep=large]
&
B_1\times_{t_B,B_0,r}Y
\arrow[dl,"\pi_Y"']
\arrow[dr,"\sigma"]
\arrow[d,"\pi_{B_1}" description]
&
\\
Y
\arrow[dr,"r"']
&
B_1
\arrow[d,"{(t_B,s_B)}" description]
&
Y
\arrow[dl,"r"]
\\
&
B_0\times B_0
&
\end{tikzcd}
.
\]
Thus \((\alpha_y,z)\) is viewed as a right-action arrow from the object over
\(q_Q(y)\) to the object over \(q_P(\theta y)\).
\end{definition}

\begin{theorem}[Naturality of the evaluated output comparison]
\label{thm:relative-output-comparison-natural}
The family
\[
\lambda_{\Theta,Y}
\]
of Definition~\ref{def:relative-downstream-comparison-transformation} is an
internal natural transformation
\[
\Lambda_{\Psi,Y}
\Rightarrow
\Lambda_{\Phi,Y}H_{\Theta,Y}.
\]
\end{theorem}

\begin{proof}
\leavevmode
\begin{enumerate}
\item \textbf{Check the components.}
The component at \((y,z)\) is well-defined because
\[
t_B\alpha_y=q_Q(y)=r(z).
\]
Its source is \(z\), and its target is
\[
\sigma(\alpha_y,z),
\]
which is the object part of
\[
\Lambda_{\Phi,Y}H_{\Theta,Y}(y,z).
\]

\item \textbf{Write the naturality square.}
Let
\[
(a,y,z):(y,z)\to
\bigl(\delta_Q(a,y),\sigma(\omega_Q(a,y),z)\bigr)
\]
be a primitive arrow of \(\mathsf{Prog}_{\Psi,Y}\). Naturality in
\(\int_{\mathbb B}Y\) asks for the commutativity of
\[
\begin{tikzcd}[column sep=large,row sep=large]
z
  \arrow[r,"{(\omega_Q(a,y),z)}"]
  \arrow[d,"{(\alpha_y,z)}"']
&
\sigma(\omega_Q(a,y),z)
  \arrow[d,"{(\alpha_{\delta_Q(a,y)},\,\sigma(\omega_Q(a,y),z))}"]
\\
\sigma(\alpha_y,z)
  \arrow[r,"{(\omega_P(a,\theta y),\,\sigma(\alpha_y,z))}"']
&
\sigma(c_{\Theta,a,y},z)
\end{tikzcd}
\]
where
\[
c_{\Theta,a,y}
:=
\omega_Q(a,y)\circ \alpha_{\delta_Q(a,y)}
=
\alpha_y\circ \omega_P(a,\theta y).
\]

\item \textbf{Make the composition order explicit.}
With the convention
\[
m_B(\beta,\gamma)=\gamma\circ\beta,
\]
the top-right composite corresponds to
\[
m_B\bigl(\alpha_{\delta_Q(a,y)},\omega_Q(a,y)\bigr)
=
\omega_Q(a,y)\circ\alpha_{\delta_Q(a,y)}.
\]
The left-bottom composite corresponds to
\[
m_B\bigl(\omega_P(a,\theta y),\alpha_y\bigr)
=
\alpha_y\circ\omega_P(a,\theta y).
\]

\item \textbf{Use the simulation law.}
The equality of these two composites is exactly the output equation in the
Mealy simulation law:
\[
\alpha_y\circ \omega_P(a,\theta y)
=
\omega_Q(a,y)\circ \alpha_{\delta_Q(a,y)}.
\]
\end{enumerate}
\end{proof}

\begin{corollary}[Relative labelled functor induced by a simulation]
\label{cor:relative-labelled-functor-simulation}
Every Mealy simulation
\[
\Theta:\Phi\Rightarrow\Psi
\]
and every right \(\mathbb B\)-action \(Y\) induce an input-strict,
downstream-lax labelled functor
\[
H_{\Theta,Y}:
(\mathsf{Prog}_{\Psi,Y},I_{\Psi,Y},\Lambda_{\Psi,Y})
\longrightarrow
(\mathsf{Prog}_{\Phi,Y},I_{\Phi,Y},\Lambda_{\Phi,Y})
\]
consisting of
\[
I_{\Phi,Y}H_{\Theta,Y}=I_{\Psi,Y}
\]
and a canonical natural transformation
\[
\lambda_{\Theta,Y}:
\Lambda_{\Psi,Y}
\Rightarrow
\Lambda_{\Phi,Y}H_{\Theta,Y}.
\]
\end{corollary}


\section{Polynomial preliminaries and dependent operators}
\label{sec:relative-polynomial-preliminaries}

The operational constructions above require only finite limits.  We now
assume the specific dependent products used below exist.  A sufficient global
hypothesis is that \(\mathcal E\) is locally cartesian closed.

\subsection{Dependent sums, pullbacks, and products}
\label{subsec:poly-relative-dependent-operators}

For a morphism
\[
f:X\to Y,
\]
write
\[
f^\ast:\mathcal E/Y\to\mathcal E/X
\]
for pullback along \(f\),
\[
\Sigma_f:\mathcal E/X\to\mathcal E/Y
\]
for postcomposition with \(f\), and
\[
\Pi_f:\mathcal E/X\to\mathcal E/Y
\]
for the right adjoint to \(f^\ast\), when it exists:
\[
\Sigma_f\dashv f^\ast\dashv\Pi_f.
\]
The adjunctions are summarized by
\[
\begin{tikzcd}[column sep=large]
\mathcal E/X
  \arrow[r,shift left=1.1ex,"\Sigma_f"]
&
\mathcal E/Y
  \arrow[l,shift left=.2ex,"f^\ast"]
  \arrow[l,shift right=1.1ex,"\Pi_f"']
\end{tikzcd}
\]
where
\[
\Sigma_f\dashv f^\ast
\qquad\text{and}\qquad
f^\ast\dashv\Pi_f.
\]

All equalities of polynomial and span constructions are read after the
canonical associativity, unit, distributivity, and Beck--Chevalley
identifications.

\section{Polynomial spans and linear spans}
\label{sec:poly-relative-polynomial-spans}

\subsection{Polynomial spans}
\label{subsec:poly-relative-polynomial-spans}

\begin{definition}[Polynomial span]
\label{def:poly-relative-polynomial-span}
A \textbf{polynomial span} from \(I\) to \(J\) is a diagram
\[
\begin{tikzcd}[column sep=large]
I
&
E
  \arrow[l,"s"']
  \arrow[r,"p"]
&
B
  \arrow[r,"t"]
&
J.
\end{tikzcd}
\]
When the required dependent product exists, it induces a functor
\[
\Sigma_t\Pi_p s^\ast:
\mathcal E/I\to\mathcal E/J.
\]
\end{definition}

Thus a polynomial span is read as the composite
\[
\begin{tikzcd}[column sep=huge]
\mathcal E/I
  \arrow[r,"s^\ast"]
&
\mathcal E/E
  \arrow[r,"\Pi_p"]
&
\mathcal E/B
  \arrow[r,"\Sigma_t"]
&
\mathcal E/J .
\end{tikzcd}
\]
This is the dependent polynomial functor associated to the displayed diagram
\cite{GambinoKock2013Polynomial}.

Fibrewise, if \(X\to I\), then
\[
\left(\Sigma_t\Pi_p s^\ast X\right)_j
\cong
\sum_{b\in B_j}
\prod_{e\in E_b}
X_{s(e)}.
\]

\subsection{Linear spans}
\label{subsec:poly-relative-linear-spans}

\begin{remark}[Linear spans]
\label{rem:poly-relative-linear-spans}
An ordinary span
\[
I\xleftarrow{s}R\xrightarrow{t}J
\]
is the linear polynomial span
\[
\begin{tikzcd}[column sep=large]
I
&
R
  \arrow[l,"s"']
  \arrow[r,"1_R"]
&
R
  \arrow[r,"t"]
&
J.
\end{tikzcd}
\]
The induced functor is
\[
\Sigma_t\Pi_{1_R}s^\ast
\cong
\Sigma_t s^\ast.
\]
Thus ordinary spans are the special case in which each execution witness has a
single continuation position.  The response polynomials below replace this
single position by a dependent family of admissible input-query positions.
\end{remark}

The polynomial-functor reading of a span is covariant:
\[
\Sigma_t s^\ast:\mathcal E/I\to\mathcal E/J.
\]
The dynamic-logic reading used later is the predicate-transformer reading:
\[
\langle R\rangle=\exists_s t^\ast.
\]
The two readings use the same span shape but act on different indexed
categories and with different variance.


\section{Response transposition of right actions}
\label{sec:relative-response-transposition-actions}

The first polynomial construction curries an ordinary right action before
output labels are added.  It provides the model for the relative response
polynomial constructed later.

\subsection{The action-response polynomial}
\label{subsec:poly-relative-action-response-polynomial}

Assume that the dependent product
\[
\Pi_{t_A}
\]
exists.

\begin{definition}[Action-response polynomial]
\label{def:poly-relative-action-response-polynomial}
The \textbf{action-response polynomial} of \(\mathbb A\) is the endofunctor on
\(\mathcal E/A_0\)
\[
\mathcal R_{\mathbb A}(X)
:=
\Pi_{t_A}s_A^\ast X.
\]
Fibrewise,
\[
\mathcal R_{\mathbb A}(X)_v
\cong
\prod_{a:u\to v}X_u.
\]
\end{definition}

The action map transposes along
\[
t_A^\ast\dashv\Pi_{t_A}
\]
to a map over \(A_0\)
\[
\widehat\rho:
X\to\mathcal R_{\mathbb A}(X).
\]
The transpose relation is displayed by
\[
\begin{tikzcd}[column sep=large,row sep=large]
t_A^\ast X
  \arrow[rr,"\rho"]
  \arrow[d]
&&
s_A^\ast X
  \arrow[d]
\\
A_1
  \arrow[rr,equal]
&&
A_1
\end{tikzcd}
\qquad\Longleftrightarrow\qquad
\begin{tikzcd}[column sep=large,row sep=large]
X
  \arrow[rr,"\widehat\rho"]
  \arrow[dr,"p"']
&&
\Pi_{t_A}s_A^\ast X
  \arrow[dl]
\\
&
A_0
&
\end{tikzcd}
\]
and in generalized-element notation
\[
\widehat\rho(x)(a)=\rho(a,x).
\]

\subsection{Multiplicativity}
\label{subsec:poly-relative-action-response-multiplicativity}

\begin{definition}[Multiplicative action-response coalgebra]
\label{def:poly-relative-multiplicative-action-coalgebra}
A coalgebra
\[
c:X\to\mathcal R_{\mathbb A}(X)
\]
is \textbf{multiplicative} if, writing
\[
c(x)(a)=x_a,
\]
it satisfies
\[
c(x)(e_A(px))=x,
\]
and, for
\[
a:u\to v,
\qquad
b:v\to w,
\qquad
p(x)=w,
\]
if
\[
c(x)(b)=x_b,
\qquad
c(x_b)(a)=x_{ab},
\]
then
\[
c(x)(m_A(a,b))=x_{ab}.
\]
\end{definition}

Here
\[
m_A(a,b)=b\circ a.
\]
Operationally, the right action executes \(b\) first at the current state and
then executes \(a\) at the updated state.

\begin{proposition}[Actions as multiplicative response coalgebras]
\label{prop:poly-relative-actions-as-response-coalgebras}
Assume \(\Pi_{t_A}\) exists.  To give a right \(\mathbb A\)-action
\[
(X\xrightarrow{p}A_0,\rho)
\]
is equivalently to give a multiplicative coalgebra
\[
c:X\to\mathcal R_{\mathbb A}(X).
\]
\end{proposition}

\begin{proof}
\leavevmode
\begin{enumerate}
\item \textbf{Transpose the action map.}
The adjunction
\[
t_A^\ast\dashv\Pi_{t_A}
\]
identifies maps
\[
t_A^\ast X\to s_A^\ast X
\]
over \(A_1\) with maps
\[
X\to\Pi_{t_A}s_A^\ast X
\]
over \(A_0\).

\item \textbf{Read the transpose in elements.}
Under this identification,
\[
\rho(a,x)
\]
corresponds to
\[
c(x)(a).
\]

\item \textbf{Compare unit laws.}
The action unit law
\[
\rho(e_A(px),x)=x
\]
is exactly
\[
c(x)(e_A(px))=x.
\]

\item \textbf{Compare multiplication laws.}
For
\[
a:u\to v,
\qquad
b:v\to w,
\qquad
p(x)=w,
\]
the action multiplication law
\[
\rho(m_A(a,b),x)=\rho(a,\rho(b,x))
\]
is exactly the multiplicativity equation
\[
c(x)(m_A(a,b))=c(c(x)(b))(a).
\]
\end{enumerate}
\end{proof}

\section{Labelled relative actions}
\label{sec:poly-relative-labelled-actions}

\subsection{Labelled actions over a downstream action}
\label{subsec:poly-relative-labelled-actions-definition}

Fix an internal category \(\mathbb B\) and a right \(\mathbb B\)-action
\[
Y=(Y\xrightarrow{r}B_0,\sigma).
\]
The action category of \(Y\) has arrow object
\[
B_1\times_{t_B,B_0,r}Y
\]
and source-target span
\[
\begin{tikzcd}
&
B_1\times_{t_B,B_0,r}Y
  \arrow[dl,"\pi_Y"']
  \arrow[dr,"\sigma"]
&
\\
Y && Y .
\end{tikzcd}
\]

\begin{definition}[Labelled relative action]
\label{def:poly-relative-labelled-action}
A \textbf{labelled relative action over \(Y\)} consists of:
\begin{enumerate}
\item a right \(\mathbb A\)-action
\[
(X\xrightarrow{p_X}A_0,\rho);
\]
\item a map
\[
y_X:X\to Y;
\]
\item a map
\[
\lambda:
A_1\times_{t_A,A_0,p_X}X\to B_1
\]
such that
\[
t_B\lambda(a,z)=r(y_Xz),
\]
and
\[
y_X(\rho(a,z))
=
\sigma(\lambda(a,z),y_Xz);
\]
\item the unit and multiplication laws saying that
\[
(a,z)\longmapsto(\lambda(a,z),y_Xz)
\]
is the arrow map of an internal functor
\[
\Lambda_X:\int_{\mathbb A}X\to\int_{\mathbb B}Y.
\]
\end{enumerate}
\end{definition}

The arrow map into the downstream action category is typed by the diagram
\[
\begin{tikzcd}[column sep=large,row sep=large]
A_1\times_{t_A,A_0,p_X}X
  \arrow[rr,"{(a,z)\mapsto(\lambda(a,z),y_Xz)}"]
  \arrow[dr,"y_X\pi_2"']
&&
B_1\times_{t_B,B_0,r}Y
  \arrow[dl,"\pi_Y"]
  \arrow[d,"\pi_{B_1}" description]
\\
&
Y
  \arrow[d,"r"']
&
B_1
  \arrow[dl,"t_B"]
  \arrow[dr,"s_B"]
\\
&
B_0
&&
B_0 .
\end{tikzcd}
\]
The left triangle records the source condition
\[
t_B\lambda=r y_X\pi_2.
\]
The target preservation condition is the triangle
\[
\begin{tikzcd}[column sep=large,row sep=large]
A_1\times_{t_A,A_0,p_X}X
  \arrow[rr,"\rho"]
  \arrow[dr,"{(\lambda,y_X\pi_2)}"']
&&
X
  \arrow[dl,"y_X"]
\\
&
Y
&
\end{tikzcd}
\]
where the diagonal map lands in the downstream arrow object and is followed by
the downstream target map \(\sigma\).

Thus a labelled relative action is equivalently a right \(\mathbb A\)-action
whose action category maps to the downstream action category
\[
\int_{\mathbb B}Y.
\]

\subsection{The labelled action associated to a Mealy morphism}
\label{subsec:poly-relative-mealy-labelled-action}

For the labelled action associated to a Mealy morphism
\[
\Phi=(P,\delta_P,\omega_P):\mathbb A\rightsquigarrow\mathbb B,
\]
one has
\[
X=S_{\Phi,Y}=P\times_{B_0}Y,
\]
\[
y_X=\pi_Y:S_{\Phi,Y}\to Y,
\]
\[
\rho(a,x,y)
=
\bigl(
\delta_P(a,x),
\sigma(\omega_P(a,x),y)
\bigr),
\]
and
\[
\lambda(a,x,y)=\omega_P(a,x).
\]
The construction is displayed by
\[
\begin{tikzcd}[column sep=large,row sep=large]
A_1\times_{t_A,A_0,p_P\pi_P}S_{\Phi,Y}
  \arrow[rr,"{\rho=\alpha_{\Phi,Y}}"]
  \arrow[dr,"{(\omega_P,\pi_Y)}"']
&&
S_{\Phi,Y}
  \arrow[dl,"\pi_Y"]
\\
&
Y
&
\end{tikzcd}
\]
where the diagonal map represents the downstream arrow
\[
(a,x,y)\longmapsto(\omega_P(a,x),y).
\]

\begin{remark}
\label{rem:poly-relative-labelled-not-bare-action}
The bare action \(\Phi^\ast Y\) remembers the next coupled state
\[
\bigl(
\delta_P(a,x),
\sigma(\omega_P(a,x),y)
\bigr).
\]
It need not remember the raw output arrow \(\omega_P(a,x)\), because different
\(\mathbb B\)-arrows may induce the same transition in \(Y\).  Output-sensitive
constructions therefore use the labelled relative action
\[
(\Phi^\ast Y,\Lambda_{\Phi,Y}),
\]
not only the underlying right \(\mathbb A\)-action.
\end{remark}

\section{The relative response interface}
\label{sec:poly-relative-interface}

\subsection{The interface object}
\label{subsec:poly-relative-interface-object}

Fix the downstream action
\[
Y=(Y\xrightarrow{r}B_0,\sigma).
\]

\begin{definition}[Relative response interface]
\label{def:poly-relative-response-interface}
The \textbf{relative response interface} is
\[
J_Y:=A_0\times Y.
\]
Write
\[
\pi_A:J_Y\to A_0,
\qquad
\pi_Y:J_Y\to Y
\]
for its projections.
\end{definition}

The product is displayed by
\[
\begin{tikzcd}
J_Y=A_0\times Y
  \arrow[r,"\pi_Y"]
  \arrow[d,"\pi_A"']
&
Y
  \arrow[d]
\\
A_0
  \arrow[r]
&
1
\arrow[ul,phantom,very near start,"\lrcorner"] .
\end{tikzcd}
\]

For a labelled relative action \(X\), define
\[
j_X:X\to J_Y,
\qquad
j_X(z)=(p_Xz,y_Xz).
\]
The defining diagram is
\[
\begin{tikzcd}[column sep=large,row sep=large]
&
X
  \arrow[dl,"p_X"']
  \arrow[dr,"y_X"]
  \arrow[d,dashed,"j_X" description]
&
\\
A_0
&
J_Y
  \arrow[l,"\pi_A"]
  \arrow[r,"\pi_Y"']
&
Y .
\end{tikzcd}
\]
For a Mealy morphism \(\Phi\), this gives
\[
j_{\Phi,Y}:S_{\Phi,Y}\to J_Y,
\qquad
j_{\Phi,Y}(x,y)=(p_Px,y).
\]

\subsection{Interfaces and simulations}
\label{subsec:poly-relative-interface-simulations}

\begin{remark}[Interface and simulations]
\label{rem:poly-relative-interface-and-simulations}
The interface \(J_Y=A_0\times Y\) is the base for strict response profiles.  A
Mealy simulation does not generally preserve this base strictly.  If
\[
\Theta:\Phi\Rightarrow\Psi
\]
has output-comparison component \(\alpha\), then
\[
h_{\Theta,Y}(y,z)
=
(\theta y,\sigma(\alpha_y,z))
\]
preserves the \(A_0\)-component, but changes the \(Y\)-component from \(z\) to
\[
\sigma(\alpha_y,z).
\]
This is why simulations become output-lax profile comparisons rather than
ordinary coalgebra morphisms over \(J_Y\).
\end{remark}

The failure of strictness over \(J_Y\) is displayed by
\[
\begin{tikzcd}[column sep=large,row sep=large]
S_{\Psi,Y}
  \arrow[r,"h_{\Theta,Y}"]
  \arrow[d,"j_{\Psi,Y}"']
&
S_{\Phi,Y}
  \arrow[d,"j_{\Phi,Y}"]
\\
J_Y
  \arrow[r,dashed,"?"]
&
J_Y .
\end{tikzcd}
\]
The \(A_0\)-component is preserved, but the \(Y\)-component is transformed by
the evaluated comparison arrow.

\section{Queries, responses, and policies}
\label{sec:poly-relative-queries-responses}

The following construction has the usual container shape: one first forms an
object of admissible input positions, then an object of possible responses at
each position, and finally a dependent product of responses over all positions
at a given interface \cite{AbbottAltenkirchGhani2005Containers}.

\subsection{Input queries}
\label{subsec:poly-relative-input-queries}

\begin{definition}[Input-query object]
\label{def:poly-relative-input-query-object}
The \textbf{input-query object} is the pullback
\[
\mathrm{Inp}_Y:=A_1\times_{t_A,A_0,\pi_A}J_Y.
\]
Thus
\[
\begin{tikzcd}[column sep=large,row sep=large]
\mathrm{Inp}_Y
  \arrow[r,"\overline a"]
  \arrow[d,"\kappa_Y"']
&
A_1
  \arrow[d,"t_A"]
\\
J_Y
  \arrow[r,"\pi_A"']
&
A_0
\arrow[ul,phantom,very near start,"\lrcorner"]
\end{tikzcd}
\]
is a pullback square.
\end{definition}

A generalized element of \(\mathrm{Inp}_Y\) is written
\[
(a,y),
\]
where \(a:u\to v\) is an arrow of \(\mathbb A\) and \(y\in Y\).  Its current
interface is
\[
\kappa_Y(a,y)=(t_Aa,y).
\]
There is also a next-input-interface map
\[
\lambda_A:\mathrm{Inp}_Y\to A_0,
\qquad
\lambda_A(a,y)=s_A(a).
\]
The maps are displayed by
\[
\begin{tikzcd}[column sep=large,row sep=large]
&
\mathrm{Inp}_Y
  \arrow[dl,"\lambda_A"']
  \arrow[d,"\kappa_Y"]
  \arrow[dr,"\overline a"]
&
\\
A_0
&
J_Y
  \arrow[l,"\pi_A"]
  \arrow[r,"\pi_Y"']
&
A_1
  \arrow[l,bend left=12,"t_A"]
  \arrow[ll,bend right=16,"s_A"']
\end{tikzcd}
\]

\subsection{Output responses}
\label{subsec:poly-relative-output-responses}

A response to a query \((a:u\to v,y)\) is an arrow of \(\mathbb B\)
\[
\beta:b'\to r(y).
\]
This arrow may be consumed by the right \(\mathbb B\)-action \(Y\), giving the
downstream continuation state
\[
\sigma(\beta,y).
\]

\begin{definition}[Output-response object]
\label{def:poly-relative-output-response-object}
The \textbf{output-response object} is
\[
\mathrm{Resp}_Y
:=
\mathrm{Inp}_Y\times_{r\pi_Y\kappa_Y,B_0,t_B}B_1.
\]
Thus
\[
\begin{tikzcd}[column sep=large,row sep=large]
\mathrm{Resp}_Y
  \arrow[r,"\overline\beta"]
  \arrow[d,"q_Y"']
&
B_1
  \arrow[d,"t_B"]
\\
\mathrm{Inp}_Y
  \arrow[r,"r\pi_Y\kappa_Y"']
&
B_0
\arrow[ul,phantom,very near start,"\lrcorner"]
\end{tikzcd}
\]
is a pullback square.
\end{definition}

A generalized element is written
\[
(a,y,\beta),
\]
where
\[
t_B(\beta)=r(y).
\]

Define the continuation-interface map
\[
n_Y:\mathrm{Resp}_Y\to J_Y
\]
by
\[
n_Y(a,y,\beta)
=
\bigl(
s_A(a),
\sigma(\beta,y)
\bigr).
\]
Equivalently,
\[
\pi_A n_Y=\lambda_A q_Y,
\qquad
\pi_Y n_Y=\sigma(\overline\beta,\pi_Y\kappa_Yq_Y).
\]
The typing of \(n_Y\) is displayed by
\[
\begin{tikzcd}[column sep=large,row sep=large]
\mathrm{Resp}_Y
  \arrow[rr,"n_Y"]
  \arrow[drr,bend left=16,"\lambda_Aq_Y"]
  \arrow[dr,"\overline\beta"']
&&
J_Y
  \arrow[d,"\pi_Y"]
  \arrow[ddl,bend left=16,"\pi_A"]
\\
&
B_1
&
Y
\\
&
A_0
&
\end{tikzcd}
\]
where the \(Y\)-component is obtained by applying the action map \(\sigma\).

\subsection{Total response policies}
\label{subsec:poly-relative-total-policies}

A total output policy at an interface
\[
(v,y)\in J_Y
\]
assigns to every input arrow
\[
a:u\to v
\]
an output arrow
\[
\beta_a:b'_a\to r(y).
\]

\begin{definition}[Policy object]
\label{def:poly-relative-policy-object}
Assume the dependent product below exists.  The \textbf{policy object} is
\[
K_{\mathbb A,\mathbb B,Y}
:=
\Pi_{\kappa_Y}(\mathrm{Resp}_Y\xrightarrow{q_Y}\mathrm{Inp}_Y).
\]
It is an object over \(J_Y\).  Write
\[
t_{\mathbb A,\mathbb B,Y}:
K_{\mathbb A,\mathbb B,Y}\to J_Y
\]
for its structure map.
\end{definition}

A generalized element of \(K_{\mathbb A,\mathbb B,Y}\) over \((v,y)\) is a
policy
\[
\pi_0:
\{a:u\to v\}
\to
\{\beta:b'\to r(y)\}.
\]

Now form the pullback
\[
E_{\mathbb A,\mathbb B,Y}
:=
\mathrm{Inp}_Y\times_{J_Y}K_{\mathbb A,\mathbb B,Y},
\]
where the maps to \(J_Y\) are \(\kappa_Y\) and
\(t_{\mathbb A,\mathbb B,Y}\).  Thus
\[
\begin{tikzcd}[column sep=large,row sep=large]
E_{\mathbb A,\mathbb B,Y}
  \arrow[r,"\epsilon_Q"]
  \arrow[d,"p_{\mathbb A,\mathbb B,Y}"']
&
\mathrm{Inp}_Y
  \arrow[d,"\kappa_Y"]
\\
K_{\mathbb A,\mathbb B,Y}
  \arrow[r,"t_{\mathbb A,\mathbb B,Y}"']
&
J_Y
\arrow[ul,phantom,very near start,"\lrcorner"]
\end{tikzcd}
\]
is a pullback square.

The counit of
\[
\kappa_Y^\ast\dashv\Pi_{\kappa_Y}
\]
gives an evaluation map
\[
\operatorname{ev}:
E_{\mathbb A,\mathbb B,Y}\to \mathrm{Resp}_Y
\]
over \(\mathrm{Inp}_Y\).  The evaluation diagram is
\[
\begin{tikzcd}[column sep=large,row sep=large]
E_{\mathbb A,\mathbb B,Y}
  \arrow[rr,"\operatorname{ev}"]
  \arrow[dr,"\epsilon_Q"']
&&
\mathrm{Resp}_Y
  \arrow[dl,"q_Y"]
\\
&
\mathrm{Inp}_Y
&
\end{tikzcd}
\]
and in generalized-element notation,
\[
\operatorname{ev}(a,y,\pi_0)=(a,y,\pi_0(a)).
\]

Define
\[
s_{\mathbb A,\mathbb B,Y}
:=
n_Y\operatorname{ev}:
E_{\mathbb A,\mathbb B,Y}\to J_Y.
\]
The structure maps are summarized by
\[
\begin{tikzcd}[column sep=large,row sep=large]
&
E_{\mathbb A,\mathbb B,Y}
  \arrow[dl,"s_{\mathbb A,\mathbb B,Y}"']
  \arrow[d,"p_{\mathbb A,\mathbb B,Y}"]
  \arrow[dr,"\operatorname{ev}"]
&
\\
J_Y
&
K_{\mathbb A,\mathbb B,Y}
  \arrow[l,"t_{\mathbb A,\mathbb B,Y}"]
&
\mathrm{Resp}_Y
  \arrow[ll,"n_Y"]
\end{tikzcd}
\]

\subsection{Dependent products required}
\label{subsec:poly-relative-dependent-products-required}

We also assume that the projection
\[
p_{\mathbb A,\mathbb B,Y}:
E_{\mathbb A,\mathbb B,Y}\to K_{\mathbb A,\mathbb B,Y}
\]
admits a dependent product.  This is automatic if \(\mathcal E\) is locally
cartesian closed.  More generally, it follows whenever \(\kappa_Y\) is
exponentiable and exponentiable morphisms are stable under pullback, since
\(p_{\mathbb A,\mathbb B,Y}\) is the pullback of \(\kappa_Y\) along
\[
t_{\mathbb A,\mathbb B,Y}:K_{\mathbb A,\mathbb B,Y}\to J_Y.
\]

\section{The relative Mealy response polynomial}
\label{sec:poly-relative-response-polynomial}

\subsection{Definition of the polynomial}
\label{subsec:poly-relative-response-polynomial-definition}

\begin{definition}[Relative Mealy response polynomial]
\label{def:poly-relative-mealy-response-polynomial}
The \textbf{relative Mealy response polynomial} associated to
\[
(\mathbb A,\mathbb B,Y)
\]
is the polynomial span
\[
\begin{tikzcd}[column sep=large]
J_Y
&
E_{\mathbb A,\mathbb B,Y}
  \arrow[l,"s_{\mathbb A,\mathbb B,Y}"']
  \arrow[r,"p_{\mathbb A,\mathbb B,Y}"]
&
K_{\mathbb A,\mathbb B,Y}
  \arrow[r,"t_{\mathbb A,\mathbb B,Y}"]
&
J_Y.
\end{tikzcd}
\]
It induces an endofunctor
\[
\mathcal M_{\mathbb A,\mathbb B,Y}:
\mathcal E/J_Y\to\mathcal E/J_Y
\]
defined by
\[
\mathcal M_{\mathbb A,\mathbb B,Y}(X)
=
\Sigma_{t_{\mathbb A,\mathbb B,Y}}
\Pi_{p_{\mathbb A,\mathbb B,Y}}
s_{\mathbb A,\mathbb B,Y}^{\ast}X.
\]
\end{definition}

In the following diagram,
\[
J_Y^{\mathrm{cur}}
\qquad\text{and}\qquad
J_Y^{\mathrm{next}}
\]
are two displayed copies of the same object \(J_Y\), used only to distinguish
the current response interface from the continuation response interface.

The full construction can be displayed as
\[
\begin{tikzcd}[column sep=large,row sep=large]
&
\mathrm{Resp}_Y
  \arrow[rr,"\overline\beta"]
  \arrow[d,"q_Y"']
  \arrow[ddd,bend left=18,"n_Y"]
&&
B_1
  \arrow[d,"t_B"]
\\
E_{\mathbb A,\mathbb B,Y}
  \arrow[ur,"\operatorname{ev}"]
  \arrow[r,"\epsilon_Q" description]
  \arrow[d,"p_{\mathbb A,\mathbb B,Y}"']
  \arrow[ddr,bend right=16,"s_{\mathbb A,\mathbb B,Y}=n_Y\operatorname{ev}"']
  \arrow[dr,phantom,very near start,"\lrcorner"]
&
\mathrm{Inp}_Y
  \arrow[rr,"r\pi_Y\kappa_Y" description]
  \arrow[d,"\kappa_Y"]
  \arrow[dr,"\overline a" description]
&&
B_0
\\
K_{\mathbb A,\mathbb B,Y}
  \arrow[r,"t_{\mathbb A,\mathbb B,Y}"']
&
J_Y^{\mathrm{cur}}
&
A_1
&
\\
&
J_Y^{\mathrm{next}}
&
&
\end{tikzcd}
\]

\subsection{Dependent distributivity}
\label{subsec:poly-relative-dependent-distributivity}

\begin{lemma}[Dependent distributivity for response policies]
\label{lem:poly-relative-dependent-distributivity}
In the locally cartesian closed setting, the dependent distributivity
isomorphism for the polynomial data gives a canonical fibrewise isomorphism
\[
\sum_{\pi_0\in\prod_a O_a}
\prod_a X_{a,\pi_0(a)}
\cong
\prod_a
\sum_{\beta\in O_a}X_{a,\beta}.
\]
Equivalently, this is the internal distributivity of dependent sums over
dependent products associated to the family of response objects over the input
query object.  It is not an external choice principle.
\end{lemma}

\begin{proof}
\leavevmode
\begin{enumerate}
\item \textbf{Define the forward map.}
Given a pair
\[
(\pi_0,(x_a)_a)
\]
where
\[
\pi_0(a)\in O_a
\qquad\text{and}\qquad
x_a\in X_{a,\pi_0(a)},
\]
send it to the family
\[
\bigl(\pi_0(a),x_a\bigr)_a
\in
\prod_a\sum_{\beta\in O_a}X_{a,\beta}.
\]

\item \textbf{Define the inverse map.}
Given a family
\[
(\beta_a,x_a)_a
\]
with
\[
\beta_a\in O_a
\qquad\text{and}\qquad
x_a\in X_{a,\beta_a},
\]
define
\[
\pi_0(a):=\beta_a.
\]
Then the second components form a family
\[
x_a\in X_{a,\pi_0(a)}.
\]

\item \textbf{Check the two composites.}
The two composites are the identity by the uniqueness clauses for dependent
sums and dependent products.
\end{enumerate}
\end{proof}

\subsection{Fibre calculation}
\label{subsec:poly-relative-fibre-calculation}

\begin{proposition}[Fibre calculation]
\label{prop:poly-relative-fibre-calculation}
Let
\[
X\to J_Y=A_0\times Y
\]
be an object over the response interface.  For every generalized interface
\[
(v,y)\in J_Y,
\]
there is a canonical fibrewise description
\[
\mathcal M_{\mathbb A,\mathbb B,Y}(X)_{(v,y)}
\cong
\prod_{a:u\to v}
\sum_{\beta:b'\to r(y)}
X_{(u,\sigma(\beta,y))}.
\]
\end{proposition}

\begin{proof}
\leavevmode
\begin{enumerate}
\item \textbf{Describe policies.}
A point of \(K_{\mathbb A,\mathbb B,Y}\) over \((v,y)\) is a total policy
\[
\pi_0:a\mapsto\beta_a
\]
assigning to each input arrow \(a:u\to v\) an output arrow
\[
\beta_a:b'_a\to r(y).
\]

\item \textbf{Describe positions over a policy.}
The object \(E_{\mathbb A,\mathbb B,Y}\) has, over such a policy, one position
for each input arrow
\[
a:u\to v.
\]
The continuation-interface map sends that position to
\[
(u,\sigma(\beta_a,y)).
\]

\item \textbf{Apply the dependent product.}
Thus
\[
\Pi_{p_{\mathbb A,\mathbb B,Y}}
s_{\mathbb A,\mathbb B,Y}^{\ast}X
\]
assigns to a policy \(\pi_0\) a family
\[
x_a\in X_{(u,\sigma(\pi_0(a),y))}
\]
for every input arrow \(a:u\to v\).

\item \textbf{Apply the dependent sum.}
Applying
\[
\Sigma_{t_{\mathbb A,\mathbb B,Y}}
\]
regards the result as lying over \((v,y)\).  Therefore the fibre is first
\[
\sum_{\pi_0}
\prod_{a:u\to v}
X_{(u,\sigma(\pi_0(a),y))}.
\]

\item \textbf{Use dependent distributivity.}
The dependent distributivity isomorphism of
Lemma~\ref{lem:poly-relative-dependent-distributivity} gives
\[
\sum_{\pi_0}
\prod_{a:u\to v}
X_{(u,\sigma(\pi_0(a),y))}
\cong
\prod_{a:u\to v}
\sum_{\beta:b'\to r(y)}
X_{(u,\sigma(\beta,y))}.
\]
\end{enumerate}
\end{proof}

\begin{remark}
\label{rem:poly-relative-response-polynomial-meaning}
At interface \((v,y)\), a total response profile must respond to every
admissible input arrow \(a:u\to v\).  For each such input, it chooses an output
arrow
\[
\beta:b'\to r(y)
\]
and a continuation state over
\[
(u,\sigma(\beta,y)).
\]
\end{remark}

\section{Coalgebras and labelled one-step data}
\label{sec:poly-relative-coalgebras}

\subsection{Coalgebras over the response interface}
\label{subsec:poly-relative-coalgebras-over-interface}

Let
\[
X\xrightarrow{j_X}J_Y
\]
be an object over the response interface.  Write
\[
p_X:=\pi_Aj_X:X\to A_0,
\qquad
y_X:=\pi_Yj_X:X\to Y.
\]
Thus
\[
\begin{tikzcd}[column sep=large,row sep=large]
&
X
  \arrow[dl,"p_X"']
  \arrow[dr,"y_X"]
  \arrow[d,"j_X" description]
&
\\
A_0
&
J_Y
  \arrow[l,"\pi_A"]
  \arrow[r,"\pi_Y"']
&
Y .
\end{tikzcd}
\]

A coalgebra
\[
c:X\to\mathcal M_{\mathbb A,\mathbb B,Y}(X)
\]
assigns to each
\[
z\in X_{(v,y)}
\]
and each input arrow
\[
a:u\to v
\]
a pair
\[
c(z)(a)=(\beta_a,z_a),
\]
where
\[
\beta_a:b'_a\to r(y),
\qquad
z_a\in X_{(u,\sigma(\beta_a,y))}.
\]

Thus \(c\) determines uncurried maps
\[
\rho_c:
A_1\times_{t_A,A_0,p_X}X\to X
\]
and
\[
\lambda_c:
A_1\times_{t_A,A_0,p_X}X\to B_1
\]
by
\[
c(z)(a)=(\lambda_c(a,z),\rho_c(a,z)).
\]

These satisfy
\[
p_X\rho_c(a,z)=s_A(a),
\]
\[
t_B\lambda_c(a,z)=r(y_Xz),
\]
and
\[
y_X\rho_c(a,z)
=
\sigma(\lambda_c(a,z),y_Xz).
\]
Conversely, such maps \((\rho,\lambda)\) determine a coalgebra by
\[
c(z)(a)=(\lambda(a,z),\rho(a,z)).
\]

\subsection{Currying theorem}
\label{subsec:poly-relative-currying-labelled-one-step}

\begin{theorem}[Currying labelled one-step data]
\label{thm:poly-relative-currying-labelled-one-step}
To give a coalgebra
\[
c:X\to\mathcal M_{\mathbb A,\mathbb B,Y}(X)
\]
is equivalently to give maps
\[
\rho:
A_1\times_{t_A,A_0,p_X}X\to X,
\qquad
\lambda:
A_1\times_{t_A,A_0,p_X}X\to B_1
\]
satisfying
\[
p_X\rho=s_A\pi_1,
\]
\[
t_B\lambda=r y_X\pi_2,
\]
and
\[
y_X\rho=\sigma(\lambda,y_X\pi_2).
\]
\end{theorem}

\begin{proof}
\leavevmode
\begin{enumerate}
\item \textbf{Start from the policy fibre.}
The polynomial was defined in policy form, so its fibre over \((v,y)\) is
initially
\[
\sum_{\pi_0\in\prod_{a:u\to v}O_a}
\prod_{a:u\to v}
X_{(u,\sigma(\pi_0(a),y))},
\]
where \(O_a\) is the fibre of possible output arrows
\[
\beta:b'\to r(y)
\]
for the input \(a:u\to v\).

\item \textbf{Distribute policies into pointwise responses.}
By Lemma~\ref{lem:poly-relative-dependent-distributivity}, this is
canonically isomorphic to
\[
\prod_{a:u\to v}
\sum_{\beta:b'\to r(y)}
X_{(u,\sigma(\beta,y))}.
\]

\item \textbf{Read a coalgebra in fibres.}
Thus a coalgebra over \(J_Y\) is equivalently a choice, for every state
\[
z\in X_{(v,y)}
\]
and every admissible input
\[
a:u\to v,
\]
of an output arrow
\[
\lambda(a,z):b'\to r(y)
\]
and a continuation state
\[
\rho(a,z)\in X_{(u,\sigma(\lambda(a,z),y))}.
\]

\item \textbf{Translate the fibre conditions into morphism equations.}
The statement that
\[
\rho(a,z)\in X_{(u,\sigma(\lambda(a,z),y))}
\]
is exactly the conjunction
\[
p_X\rho=s_A\pi_1,
\]
\[
t_B\lambda=r y_X\pi_2,
\]
and
\[
y_X\rho=\sigma(\lambda,y_X\pi_2).
\]
\end{enumerate}
\end{proof}

\section{Multiplicative labelled response coalgebras}
\label{sec:poly-relative-multiplicative-coalgebras}

\subsection{Definition of multiplicativity}
\label{subsec:poly-relative-multiplicative-definition}

The coalgebra map records typed one-step labelled response data.  To correspond
to an action category, this profile must preserve identities and composition.

We keep the composition convention of the previous chapters: for internally
composable arrows
\[
x\xrightarrow{\alpha}y\xrightarrow{\beta}z
\]
in \(\mathbb B\),
\[
m_B(\alpha,\beta)=\beta\circ\alpha.
\]

\begin{definition}[Multiplicative labelled response coalgebra]
\label{def:poly-relative-multiplicative-response-coalgebra}
A coalgebra
\[
c:X\to\mathcal M_{\mathbb A,\mathbb B,Y}(X)
\]
is \textbf{multiplicative} if, writing
\[
c(z)(a)=(\beta_a,z_a),
\]
the following equations hold.
\begin{enumerate}
\item For every \(z\in X\),
\[
c(z)(e_A(p_Xz))
=
(e_B(r(y_Xz)),z).
\]

\item For every composable pair
\[
a:u\to v,
\qquad
b:v\to w,
\]
and every \(z\in X_{(w,y)}\), if
\[
c(z)(b)=(\beta,z')
\]
and
\[
c(z')(a)=(\alpha,z'')
\]
then
\[
c(z)(m_A(a,b))
=
(m_B(\alpha,\beta),z'').
\]
\end{enumerate}
\end{definition}

Here \(m_A(a,b)=b\circ a\), but operationally the right action executes \(b\)
first at the current state and then executes \(a\) at the updated state.
Similarly, the output arrow for the composite is
\[
m_B(\alpha,\beta)=\beta\circ\alpha.
\]

\subsection{Multiplicative coalgebras as labelled actions}
\label{subsec:poly-relative-coalgebras-labelled-actions}

\begin{proposition}[Multiplicative coalgebras as labelled actions]
\label{prop:poly-relative-coalgebras-labelled-actions}
Assume the relevant dependent products exist.  To give a multiplicative
coalgebra
\[
c:X\to\mathcal M_{\mathbb A,\mathbb B,Y}(X)
\]
is equivalently to give:
\begin{enumerate}
\item a right \(\mathbb A\)-action on \(X\xrightarrow{p_X}A_0\);
\item an internal functor
\[
\Lambda_X:\int_{\mathbb A}X\to\int_{\mathbb B}Y
\]
whose object map is \(y_X:X\to Y\).
\end{enumerate}
Equivalently, it is a labelled relative action over \(Y\).
\end{proposition}

\begin{proof}
\leavevmode
\begin{enumerate}
\item \textbf{Uncurry the coalgebra.}
By Theorem~\ref{thm:poly-relative-currying-labelled-one-step}, a coalgebra is
equivalent to maps
\[
\rho:
A_1\times_{t_A,A_0,p_X}X\to X,
\qquad
\lambda:
A_1\times_{t_A,A_0,p_X}X\to B_1
\]
with the displayed typing equations.

\item \textbf{Recover the right action.}
The first component \(\rho\), together with the identity and composition
clauses of multiplicativity, is exactly a right \(\mathbb A\)-action on
\(X\).

\item \textbf{Construct the action category.}
The action category \(\int_{\mathbb A}X\) has object object \(X\) and arrow
object
\[
A_1\times_{t_A,A_0,p_X}X.
\]

\item \textbf{Construct the downstream functor.}
The map
\[
(a,z)\longmapsto(\lambda(a,z),y_Xz)
\]
is well-defined as a map into the arrow object of \(\int_{\mathbb B}Y\)
because
\[
t_B\lambda(a,z)=r(y_Xz).
\]
Its source and target compatibility are displayed by
\[
\begin{tikzcd}[column sep=large,row sep=large]
A_1\times_{t_A,A_0,p_X}X
  \arrow[rr,"{(\lambda,y_X\pi_2)}"]
  \arrow[d,"\pi_2"']
&&
B_1\times_{t_B,B_0,r}Y
  \arrow[d,"\pi_Y"]
\\
X
  \arrow[rr,"y_X"']
&&
Y
\end{tikzcd}
\]
and
\[
\begin{tikzcd}[column sep=large,row sep=large]
A_1\times_{t_A,A_0,p_X}X
  \arrow[rr,"{(\lambda,y_X\pi_2)}"]
  \arrow[d,"\rho"']
&&
B_1\times_{t_B,B_0,r}Y
  \arrow[d,"\sigma"]
\\
X
  \arrow[rr,"y_X"']
&&
Y .
\end{tikzcd}
\]

\item \textbf{Compare identity and composition.}
The unit and multiplication clauses of multiplicativity are precisely identity
and composition preservation for this internal functor.
\end{enumerate}
\end{proof}

\section{The coalgebra induced by a Mealy morphism}
\label{sec:poly-relative-induced-coalgebra}

\subsection{Definition of the induced coalgebra}
\label{subsec:poly-relative-induced-coalgebra-definition}

Let
\[
\Phi=(P,\delta_P,\omega_P):\mathbb A\rightsquigarrow\mathbb B
\]
be a Mealy morphism and let
\[
Y=(Y\xrightarrow{r}B_0,\sigma)
\]
be a right \(\mathbb B\)-action.

Recall that
\[
S_{\Phi,Y}=P\times_{B_0}Y
\]
maps to
\[
J_Y=A_0\times Y
\]
by
\[
j_{\Phi,Y}(x,y)=(p_Px,y).
\]
The structure map is displayed by
\[
\begin{tikzcd}[column sep=large,row sep=large]
&
S_{\Phi,Y}
  \arrow[dl,"p_P\pi_P"']
  \arrow[dr,"\pi_Y"]
  \arrow[d,"j_{\Phi,Y}" description]
&
\\
A_0
&
J_Y
  \arrow[l,"\pi_A"]
  \arrow[r,"\pi_Y"']
&
Y .
\end{tikzcd}
\]

\begin{definition}[Relative Mealy response coalgebra]
\label{def:poly-relative-mealy-response-coalgebra}
The \textbf{relative Mealy response coalgebra} induced by \(\Phi\) at \(Y\) is
the coalgebra
\[
c_{\Phi,Y}:
S_{\Phi,Y}\to
\mathcal M_{\mathbb A,\mathbb B,Y}(S_{\Phi,Y})
\]
given fibrewise by
\[
c_{\Phi,Y}(x,y)(a)
=
\left(
\omega_P(a,x),
\left(
\delta_P(a,x),
\sigma(\omega_P(a,x),y)
\right)
\right).
\]
\end{definition}

\subsection{Transpose of the labelled action}
\label{subsec:poly-relative-mealy-coalgebra-transpose}

\begin{theorem}[Mealy coalgebra as transpose of labelled action]
\label{thm:poly-relative-mealy-coalgebra-transpose}
The coalgebra
\[
c_{\Phi,Y}
\]
is the dependent-product transpose of the labelled relative action
\[
(\Phi^\ast Y,\Lambda_{\Phi,Y}).
\]
Moreover, it is multiplicative.
\end{theorem}

\begin{proof}
\leavevmode
\begin{enumerate}
\item \textbf{Recall the action map.}
The action map of \(\Phi^\ast Y\) is
\[
(a,x,y)
\longmapsto
\left(
\delta_P(a,x),
\sigma(\omega_P(a,x),y)
\right).
\]

\item \textbf{Recall the downstream label.}
The downstream comparison functor supplies the output label
\[
(a,x,y)\longmapsto\omega_P(a,x).
\]

\item \textbf{Combine the two components.}
Together these give the uncurried labelled one-step data
\[
(a,x,y)
\longmapsto
\left(
\omega_P(a,x),
\left(
\delta_P(a,x),
\sigma(\omega_P(a,x),y)
\right)
\right).
\]

\item \textbf{Transpose.}
By Theorem~\ref{thm:poly-relative-currying-labelled-one-step}, this
uncurried labelled map transposes to \(c_{\Phi,Y}\).

\item \textbf{Check multiplicativity.}
Multiplicativity follows from the Mealy unit and multiplication equations for
\(\Phi\) and from the unit and multiplication laws of the right
\(\mathbb B\)-action \(Y\).  Equivalently, it follows from
Proposition~\ref{prop:poly-relative-coalgebras-labelled-actions}, since
\[
\Lambda_{\Phi,Y}:\mathsf{Prog}_{\Phi,Y}\to\int_{\mathbb B}Y
\]
is an internal functor.
\end{enumerate}
\end{proof}

\section{Flattening response coalgebras}
\label{sec:poly-relative-flattening}

\subsection{Position objects}
\label{subsec:poly-relative-position-objects}

Let
\[
J_Y
\xleftarrow{s}
E_{\mathbb A,\mathbb B,Y}
\xrightarrow{p}
K_{\mathbb A,\mathbb B,Y}
\xrightarrow{t}
J_Y
\]
be the relative response polynomial.

For an object
\[
X\to J_Y,
\]
define the position object at \(X\) by
\[
\mathrm{Pos}_X
:=
E_{\mathbb A,\mathbb B,Y}
\times_{K_{\mathbb A,\mathbb B,Y}}
\Pi_p(s^\ast X).
\]
It is displayed by the pullback
\[
\begin{tikzcd}[column sep=large,row sep=large]
\mathrm{Pos}_X
  \arrow[r]
  \arrow[d]
&
\Pi_p(s^\ast X)
  \arrow[d]
\\
E_{\mathbb A,\mathbb B,Y}
  \arrow[r,"p"']
&
K_{\mathbb A,\mathbb B,Y}
\arrow[ul,phantom,very near start,"\lrcorner"] .
\end{tikzcd}
\]
It carries a canonical projection
\[
\partial_X:\mathrm{Pos}_X\to
\mathcal M_{\mathbb A,\mathbb B,Y}(X).
\]
Externally, a point of \(\mathrm{Pos}_X\) is a total response profile together
with a chosen input position.  This is the usual passage from a container-like
total object to its object of positions
\cite{AbbottAltenkirchGhani2005Containers}.

The position object has evaluation maps
\[
\operatorname{in}:\mathrm{Pos}_X\to A_1,
\]
\[
\operatorname{out}:\mathrm{Pos}_X\to B_1,
\]
\[
\operatorname{down}:\mathrm{Pos}_X\to
\left(\int_{\mathbb B}Y\right)_1,
\]
and
\[
\operatorname{next}:\mathrm{Pos}_X\to X.
\]
They extract the chosen input arrow, raw output arrow, downstream transition,
and continuation state.  These maps fit into
\[
\begin{tikzcd}[column sep=large,row sep=large]
&
\mathrm{Pos}_X
  \arrow[dl,"\operatorname{in}"']
  \arrow[d,"\operatorname{down}" description]
  \arrow[dr,"\operatorname{next}"]
&
\\
A_1
&
\left(\int_{\mathbb B}Y\right)_1
  \arrow[d,"\pi_{B_1}" description]
  \arrow[r,"\sigma"]
&
Y
\\
&
B_1
&
\end{tikzcd}
\]
with
\[
\operatorname{out}
=
\pi_{B_1}\operatorname{down},
\qquad
y_X\operatorname{next}
=
\sigma\operatorname{down}.
\]

\subsection{Flattening a coalgebra}
\label{subsec:poly-relative-flattening-coalgebra}

Let
\[
c:X\to\mathcal M_{\mathbb A,\mathbb B,Y}(X)
\]
be a coalgebra.  Pulling back the position projection along \(c\) gives
\[
\begin{tikzcd}[column sep=large,row sep=large]
M_c
  \arrow[r]
  \arrow[d,"s_c"']
&
\mathrm{Pos}_X
  \arrow[d,"\partial_X"]
\\
X
  \arrow[r,"c"']
&
\mathcal M_{\mathbb A,\mathbb B,Y}(X)
\arrow[ul,phantom,very near start,"\lrcorner"] .
\end{tikzcd}
\]
The object \(M_c\) is the flattened one-step execution object of \(c\).

Under the uncurried presentation of \(c\), one has
\[
M_c\cong A_1\times_{t_A,A_0,p_X}X.
\]
The source map is
\[
s_c(a,z)=z,
\]
and the target map is
\[
t_c(a,z)=\rho_c(a,z).
\]
The flattened execution span is
\[
\begin{tikzcd}
&
M_c
  \arrow[dl,"s_c"']
  \arrow[dr,"t_c"]
&
\\
X && X .
\end{tikzcd}
\]

\subsection{Flattening multiplicative coalgebras}
\label{subsec:poly-relative-flattening-multiplicative}

\begin{theorem}[Flattening multiplicative coalgebras]
\label{thm:poly-relative-flattening-multiplicative}
If
\[
c:X\to\mathcal M_{\mathbb A,\mathbb B,Y}(X)
\]
is multiplicative, then the flattened span
\[
X\xleftarrow{s_c}M_c\xrightarrow{t_c}X
\]
is the one-step execution span of the action category
\[
\int_{\mathbb A}X.
\]
Moreover, the evaluation maps on \(M_c\) recover the input projection and the
downstream comparison functor
\[
\int_{\mathbb A}X\to\int_{\mathbb B}Y.
\]
\end{theorem}

\begin{proof}
\leavevmode
\begin{enumerate}
\item \textbf{Replace the coalgebra by labelled action data.}
By Proposition~\ref{prop:poly-relative-coalgebras-labelled-actions}, a
multiplicative coalgebra is equivalently a right \(\mathbb A\)-action on
\(X\) equipped with an internal functor
\[
\Lambda_X:\int_{\mathbb A}X\to\int_{\mathbb B}Y.
\]

\item \textbf{Identify the flattened position object.}
The flattened position object is exactly
\[
A_1\times_{t_A,A_0,p_X}X,
\]
with source \(z\) and target \(\rho_c(a,z)\).

\item \textbf{Identify the action-category span.}
These are the source and target of the action category
\[
\int_{\mathbb A}X.
\]

\item \textbf{Read the evaluation maps.}
The evaluation maps give the chosen input arrow, output arrow, downstream
transition, and continuation state, hence recover the input projection and the
arrow map of \(\Lambda_X\).
\end{enumerate}
\end{proof}

\subsection{Flattening the Mealy coalgebra}
\label{subsec:poly-relative-flattening-mealy-coalgebra}

\begin{lemma}[Flattening pullback for the Mealy coalgebra]
\label{lem:poly-relative-mealy-position-pullback}
For the relative Mealy response coalgebra
\[
c_{\Phi,Y}:S_{\Phi,Y}\to
\mathcal M_{\mathbb A,\mathbb B,Y}(S_{\Phi,Y}),
\]
the pullback of the position projection
\[
\partial_S:\mathrm{Pos}_S\to
\mathcal M_{\mathbb A,\mathbb B,Y}(S_{\Phi,Y})
\]
along \(c_{\Phi,Y}\) is canonically isomorphic to
\[
M_{\Phi,Y}
=
A_1\times_{t_A,A_0,p_P\pi_P}S_{\Phi,Y}.
\]
Under this isomorphism,
\[
\operatorname{in}=\iota_A,
\qquad
\operatorname{out}=\iota_B,
\qquad
\operatorname{down}=\lambda_{\Phi,Y},
\qquad
\operatorname{next}=t_{\Phi,Y}.
\]
\end{lemma}

\begin{proof}
\leavevmode
\begin{enumerate}
\item \textbf{Pull back positions along the coalgebra.}
Pulling back \(\partial_S\) along \(c_{\Phi,Y}\) chooses, from the total
response profile determined by a state \((x,y)\), a single input position
\[
a:u\to p_P(x).
\]

\item \textbf{Identify the resulting object.}
Thus the resulting object is precisely the object of primitive executions
\[
A_1\times_{t_A,A_0,p_P\pi_P}S_{\Phi,Y}.
\]

\item \textbf{Identify the evaluation maps.}
The four evaluation maps read off, respectively, the input arrow, the raw
emitted output, the downstream transition, and the next coupled state.  These
are exactly
\[
\iota_A,\qquad
\iota_B,\qquad
\lambda_{\Phi,Y},\qquad
t_{\Phi,Y}.
\]
\end{enumerate}
\end{proof}

\begin{theorem}[Program category recovered by flattening]
\label{thm:poly-relative-program-category-recovered}
Flattening the relative Mealy response coalgebra
\[
c_{\Phi,Y}
\]
recovers the one-step execution span
\[
\mathbf m_{\Phi,Y}
=
\left(
S_{\Phi,Y}
\xleftarrow{s_{\Phi,Y}}
M_{\Phi,Y}
\xrightarrow{t_{\Phi,Y}}
S_{\Phi,Y}
\right)
\]
of the relative program category
\[
\mathsf{Prog}_{\Phi,Y}.
\]
Explicitly,
\[
M_{\Phi,Y}
=
A_1\times_{t_A,A_0,p_P\pi_P}S_{\Phi,Y},
\]
\[
s_{\Phi,Y}(a,x,y)=(x,y),
\]
and
\[
t_{\Phi,Y}(a,x,y)
=
\left(
\delta_P(a,x),
\sigma(\omega_P(a,x),y)
\right).
\]
Consequently, the internal category recovered from the flattened coalgebra is
\[
\mathsf{Prog}_{\Phi,Y}
=
\int_{\mathbb A}\Phi^\ast Y.
\]
\end{theorem}

\begin{proof}
\leavevmode
\begin{enumerate}
\item \textbf{Use multiplicativity.}
By Theorem~\ref{thm:poly-relative-mealy-coalgebra-transpose},
\(c_{\Phi,Y}\) is multiplicative.

\item \textbf{Apply flattening.}
Theorem~\ref{thm:poly-relative-flattening-multiplicative} identifies the
flattened span with the one-step execution span of the action category.

\item \textbf{Use the Mealy pullback calculation.}
Lemma~\ref{lem:poly-relative-mealy-position-pullback} identifies this flattened
one-step object with
\[
M_{\Phi,Y}.
\]

\item \textbf{Read the explicit formulas.}
The source and target maps are the uncurried maps defining the labelled action
associated to \(\Phi\).
\end{enumerate}
\end{proof}

\begin{remark}
\label{rem:poly-relative-polynomial-not-replacement}
The polynomial coalgebra does not replace the program category.  It is the
curried total-response presentation of the same labelled one-step data whose
flattened form is the one-step execution span.
\end{remark}

\section{Output-lax profile comparisons}
\label{sec:poly-relative-output-lax-profile-comparisons}

\subsection{Why ordinary coalgebra morphisms are too strict}
\label{subsec:poly-relative-ordinary-morphisms-too-strict}

The relative response polynomial lives over
\[
J_Y=A_0\times Y.
\]
Ordinary coalgebra morphisms over \(J_Y\) are too strict for Mealy simulations,
because simulations execute output-comparison arrows in \(Y\) and therefore
change the \(Y\)-component of the interface.

Thus we isolate the correct comparison structure.

\subsection{Definition of output-lax comparison}
\label{subsec:poly-relative-output-lax-comparison-definition}

Let
\[
c:X\to\mathcal M_{\mathbb A,\mathbb B,Y}(X),
\qquad
d:Z\to\mathcal M_{\mathbb A,\mathbb B,Y}(Z)
\]
be multiplicative labelled response coalgebras.  Write their uncurried data as
\[
(\rho_X,\lambda_X)
\qquad\text{and}\qquad
(\rho_Z,\lambda_Z).
\]

\begin{definition}[Output-lax profile comparison]
\label{def:poly-relative-output-lax-profile-comparison}
An \textbf{output-lax profile comparison}
\[
(c:X\to\mathcal M(X))
\longrightarrow
(d:Z\to\mathcal M(Z))
\]
consists of maps
\[
h:X\to Z,
\qquad
\chi:X\to B_1,
\]
such that
\[
p_Zh=p_X,
\]
\[
t_B\chi=r y_X,
\qquad
s_B\chi=r y_Zh,
\]
and
\[
y_Zh=\sigma(\chi,y_X).
\]
Thus \(\chi_x\) is an arrow
\[
r(y_Zh(x))\longrightarrow r(y_Xx)
\]
in \(\mathbb B\), and it evaluates in the right action \(Y\) as an arrow
\[
y_Xx\longrightarrow y_Zh(x)
\]
of \(\int_{\mathbb B}Y\).

Writing
\[
c(x)(a)=(\beta_X(a,x),x_a)
\]
and
\[
d(hx)(a)=(\beta_Z(a,hx),z_a),
\]
the comparison also satisfies the strict input-target equation
\[
z_a=h(x_a),
\]
or equivalently
\[
\rho_Z(a,hx)=h(\rho_X(a,x)).
\]
Finally, the output-lax naturality equation is
\[
\chi_x\circ \beta_Z(a,hx)
=
\beta_X(a,x)\circ \chi_{x_a}.
\]
Equivalently,
\[
m_B(\beta_Z(a,hx),\chi_x)
=
m_B(\chi_{x_a},\beta_X(a,x)).
\]
\end{definition}

The structure equations can be displayed as
\[
\begin{tikzcd}[column sep=large,row sep=large]
&
X
  \arrow[dl,"y_X"']
  \arrow[dr,"y_Zh"]
  \arrow[d,"\chi" description]
&
\\
Y
  \arrow[d,"r"']
&
B_1
  \arrow[dl,"t_B" description]
  \arrow[dr,"s_B" description]
&
Y
  \arrow[d,"r"]
\\
B_0
&&
B_0 .
\end{tikzcd}
\]
The action condition is the additional equation
\[
y_Zh=\sigma(\chi,y_X).
\]

\subsection{Flattening output-lax comparisons}
\label{subsec:poly-relative-flatten-output-lax-comparisons}

\begin{proposition}[Flattening an output-lax profile comparison]
\label{prop:poly-relative-flatten-output-lax-profile-comparison}
An output-lax profile comparison \((h,\chi):c\to d\) induces:
\begin{enumerate}
\item an internal functor
\[
H:\int_{\mathbb A}X\to\int_{\mathbb A}Z
\]
strictly over \(\mathbb A^{\mathrm{op}}\);
\item an internal natural transformation
\[
\lambda^\chi:
\Lambda_X\Rightarrow\Lambda_ZH
\]
with codomain \(\int_{\mathbb B}Y\).
\end{enumerate}
The functor \(H\) is given by
\[
H_0=h,
\qquad
H_1(a,x)=(a,hx),
\]
and the natural transformation has component
\[
(\chi_x,y_Xx):y_Xx\to y_Zh(x)
\]
at \(x\).
\end{proposition}

\begin{proof}
\leavevmode
\begin{enumerate}
\item \textbf{Define the object map.}
Set
\[
H_0=h:X\to Z.
\]
The equation
\[
p_Zh=p_X
\]
says that this map is over \(A_0\).

\item \textbf{Define the arrow map.}
The arrow map is
\[
H_1(a,x)=(a,hx).
\]
It is well-defined because
\[
t_A(a)=p_Xx=p_Zh(x).
\]

\item \textbf{Check source preservation.}
Source preservation is immediate from
\[
s(a,x)=x
\]
and
\[
s(a,hx)=hx.
\]

\item \textbf{Check target preservation.}
The strict input-target equation
\[
\rho_Z(a,hx)=h(\rho_X(a,x))
\]
says exactly that \(H\) preserves targets of primitive arrows.

\item \textbf{Check identities and composition.}
Since the input arrow \(a\) is unchanged and \(h\) preserves the action targets,
identity and composition preservation follow from multiplicativity of the two
coalgebras.

\item \textbf{Construct the downstream natural transformation.}
The component
\[
(\chi_x,y_Xx)
\]
is a well-defined arrow of \(\int_{\mathbb B}Y\) because
\[
t_B\chi_x=r(y_Xx)
\]
and
\[
y_Zh(x)=\sigma(\chi_x,y_Xx).
\]

\item \textbf{Check naturality.}
The output-lax naturality equation
\[
\chi_x\circ \beta_Z(a,hx)
=
\beta_X(a,x)\circ \chi_{x_a}
\]
is exactly the naturality square in the downstream action category.
\end{enumerate}
\end{proof}

The naturality square is
\[
\begin{tikzcd}[column sep=large,row sep=large]
y_Xx
  \arrow[r,"{(\beta_X(a,x),y_Xx)}"]
  \arrow[d,"{(\chi_x,y_Xx)}"']
&
y_Xx_a
  \arrow[d,"{(\chi_{x_a},y_Xx_a)}"]
\\
y_Zh(x)
  \arrow[r,"{(\beta_Z(a,hx),y_Zh(x))}"']
&
y_Zh(x_a).
\end{tikzcd}
\]

\section{Mealy simulations as output-lax profile comparisons}
\label{sec:poly-relative-mealy-simulations-profile-comparisons}

\subsection{Evaluated comparison data}
\label{subsec:poly-relative-simulation-evaluated-data}

Let
\[
\Phi=(P,\delta_P,\omega_P),
\qquad
\Psi=(Q,\delta_Q,\omega_Q)
:
\mathbb A\rightsquigarrow\mathbb B
\]
be Mealy morphisms, and let
\[
\Theta:\Phi\Rightarrow\Psi
\]
be a Mealy simulation with components
\[
\theta:Q\to P,
\qquad
\alpha:Q\to B_1.
\]
Thus
\[
p_P\theta=p_Q,
\qquad
s_B\alpha=q_P\theta,
\qquad
t_B\alpha=q_Q,
\]
and
\[
\theta(\delta_Q(a,y))
=
\delta_P(a,\theta y),
\]
\[
\alpha_y\circ\omega_P(a,\theta y)
=
\omega_Q(a,y)\circ\alpha_{\delta_Q(a,y)}.
\]

Let \(Y\) be a right \(\mathbb B\)-action.  The evaluated state comparison is
\[
h_{\Theta,Y}:S_{\Psi,Y}\to S_{\Phi,Y},
\qquad
h_{\Theta,Y}(y,z)
=
\bigl(\theta y,\sigma(\alpha_y,z)\bigr).
\]
Define
\[
\chi_{\Theta,Y}:S_{\Psi,Y}\to B_1
\]
by
\[
\chi_{\Theta,Y}(y,z)=\alpha_y.
\]

\subsection{Simulations induce output-lax comparisons}
\label{subsec:poly-relative-simulation-output-lax-theorem}

\begin{theorem}[Simulations induce output-lax profile comparisons]
\label{thm:poly-relative-simulations-profile-comparisons}
The pair
\[
(h_{\Theta,Y},\chi_{\Theta,Y})
\]
is an output-lax profile comparison
\[
c_{\Psi,Y}\longrightarrow c_{\Phi,Y}.
\]
After flattening, it recovers the input-strict functor
\[
H_{\Theta,Y}:
\mathsf{Prog}_{\Psi,Y}\to\mathsf{Prog}_{\Phi,Y}
\]
and the downstream comparison natural transformation
\[
\lambda_{\Theta,Y}:
\Lambda_{\Psi,Y}
\Rightarrow
\Lambda_{\Phi,Y}H_{\Theta,Y}.
\]
\end{theorem}

\begin{proof}
\leavevmode
\begin{enumerate}
\item \textbf{Check input strictness.}
The map \(h_{\Theta,Y}\) is strict over \(A_0\) because
\[
p_P\theta=p_Q.
\]

\item \textbf{Check the output-comparison typing.}
The map \(\chi_{\Theta,Y}\) satisfies
\[
t_B\chi_{\Theta,Y}(y,z)=t_B\alpha_y=q_Qy=r(z),
\]
and
\[
s_B\chi_{\Theta,Y}(y,z)=s_B\alpha_y=q_P\theta y
=
r(\sigma(\alpha_y,z)).
\]

\item \textbf{Check evaluation in the downstream action.}
Moreover,
\[
\pi_Yh_{\Theta,Y}(y,z)=\sigma(\alpha_y,z)
=
\sigma(\chi_{\Theta,Y}(y,z),z).
\]

\item \textbf{Compute the source response.}
The source coalgebra \(c_{\Psi,Y}\) has
\[
c_{\Psi,Y}(y,z)(a)
=
\left(
\omega_Q(a,y),
\left(
\delta_Q(a,y),
\sigma(\omega_Q(a,y),z)
\right)
\right).
\]

\item \textbf{Compute the target response at the compared state.}
The target coalgebra at \(h_{\Theta,Y}(y,z)\) has
\[
c_{\Phi,Y}(\theta y,\sigma(\alpha_y,z))(a)
=
\left(
\omega_P(a,\theta y),
\left(
\delta_P(a,\theta y),
\sigma(\omega_P(a,\theta y),\sigma(\alpha_y,z))
\right)
\right).
\]

\item \textbf{Check strict target preservation.}
The strict input-target equation is
\[
\rho_{\Phi}(a,h_{\Theta,Y}(y,z))
=
h_{\Theta,Y}(\rho_{\Psi}(a,y,z)).
\]
Indeed, the left-hand side is
\[
\left(
\delta_P(a,\theta y),
\sigma(\omega_P(a,\theta y),\sigma(\alpha_y,z))
\right),
\]
while the right-hand side is
\[
\left(
\theta\delta_Q(a,y),
\sigma(\alpha_{\delta_Q(a,y)},\sigma(\omega_Q(a,y),z))
\right).
\]
The first components agree by the state part of the simulation law.  The second
components agree by the output part of the simulation law and the right-action
multiplication law for \(Y\).

\item \textbf{Check output-lax naturality.}
The output-lax square is precisely
\[
\alpha_y\circ\omega_P(a,\theta y)
=
\omega_Q(a,y)\circ\alpha_{\delta_Q(a,y)}.
\]
Therefore \((h_{\Theta,Y},\chi_{\Theta,Y})\) is an output-lax profile
comparison.

\item \textbf{Flatten.}
The final statement is
Proposition~\ref{prop:poly-relative-flatten-output-lax-profile-comparison}
specialized to the Mealy-induced coalgebras.
\end{enumerate}
\end{proof}

\begin{remark}
\label{rem:poly-relative-not-ordinary-coalgebra-map}
The map
\[
h_{\Theta,Y}:S_{\Psi,Y}\to S_{\Phi,Y}
\]
is generally not a map over \(J_Y=A_0\times Y\), since it changes the
downstream component from \(z\) to \(\sigma(\alpha_y,z)\).  The correct
polynomial comparison is therefore output-lax, not an ordinary coalgebra
morphism in \(\mathcal E/J_Y\).
\end{remark}


\section{Finite trajectories of response coalgebras}
\label{sec:relative-response-temporalization}

A multiplicative response coalgebra carries no temporal structure beyond the
labelled action already encoded by its multiplication laws.  Its finite
trajectories are therefore obtained by flattening and applying the path
construction of the preceding chapter.

\begin{definition}[Trajectory object of a multiplicative response coalgebra]
\label{def:relative-response-coalgebra-trajectories}
Let
\[
    c:X\to\mathcal M_{\mathbb A,\mathbb B,Y}(X)
\]
be multiplicative, and let
\[
    \Lambda_X:\int_{\mathbb A}X\to\int_{\mathbb B}Y
\]
be the corresponding labelled relative action.  Define
\[
    \mathscr T_c
    :=
    \operatorname{Path}\!\left(\int_{\mathbb A}X\right).
\]
\end{definition}

\begin{proposition}[Currying does not change finite execution]
\label{prop:relative-currying-preserves-trajectories}
Under the equivalence between multiplicative response coalgebras and labelled
relative actions, the trajectory object \(\mathscr T_c\) depends only on the
uncurried labelled action.  Flattening \(c\) and then forming finite paths
recovers the same labelled trajectory object.
\end{proposition}

\begin{proof}
By Proposition~\ref{prop:poly-relative-coalgebras-labelled-actions}, the
multiplicative coalgebra and the labelled action are two transposes of the
same maps \((\rho,\lambda)\).  By
Theorem~\ref{thm:poly-relative-flattening-multiplicative}, flattening recovers
the one-step span and internal-category operations of
\(\int_{\mathbb A}X\), together with the downstream labelling functor.
Applying \(\operatorname{Path}\) therefore gives the same labelled path
object in every degree.
\end{proof}

\begin{corollary}[Relative Mealy trajectories from the response coalgebra]
\label{cor:relative-mealy-response-trajectories}
For the canonical response coalgebra \(c_{\Phi,Y}\),
\[
    \mathscr T_{c_{\Phi,Y}}
    \cong
    \operatorname{Path}(\mathsf{Prog}_{\Phi,Y})
\]
and hence
\[
    \mathscr T_{c_{\Phi,Y}}
    \cong
    \mathscr T_\Phi
    \times_{\operatorname{Path}(\mathbb B^{\mathrm{op}})}
    \operatorname{Path}\!\left(\int_{\mathbb B}Y\right).
\]
\end{corollary}

\begin{proof}
The first isomorphism follows from
Theorem~\ref{thm:poly-relative-program-category-recovered}.  The second is
Corollary~\ref{cor:relative-finite-trajectories-pullback}.
\end{proof}

\begin{remark}[Multiplicativity as curried Segal coherence]
\label{rem:relative-multiplicativity-segal}
The unit and multiplication equations of a response coalgebra are the
curried form of identity insertion and contraction in its execution category.
Thus the Segal coherence of finite execution and the multiplicativity of the
total response profile express the same categorical structure in uncurried
and curried form, respectively.
\end{remark}

\section{The one-object specialization}
\label{sec:poly-relative-one-object-specialization}

\subsection{One-object Mealy response profiles}
\label{subsec:poly-relative-one-object-response-profiles}

Let \(M\) and \(N\) be monoid objects in \(\mathcal E\), regarded as
one-object internal categories with the convention
\[
b\circ a=ba.
\]
A Mealy morphism
\[
\Phi:M\rightsquigarrow N
\]
consists of an object \(P\), a right \(M\)-action
\[
x\cdot m:=\delta_P(m,x),
\]
and an \(N\)-valued output cocycle
\[
\omega_P:M\times P\to N
\]
satisfying
\[
\omega_P(1,x)=1,
\qquad
\omega_P(mn,x)=\omega_P(m,x)\omega_P(n,x\cdot m).
\]

Let \(Y\) be a right \(N\)-object.  Then
\[
J_Y=1\times Y\cong Y.
\]
The relative response polynomial has fibre
\[
\mathcal M_{M,N,Y}(X)_y
\cong
\prod_{m\in M}
\sum_{n\in N}
X_{y\cdot n}.
\]
The induced coalgebra is
\[
c_{\Phi,Y}(x,y)(m)
=
\left(
\omega_P(m,x),
\left(
x\cdot m,
y\cdot\omega_P(m,x)
\right)
\right).
\]

\begin{proposition}[One-object cascade response coalgebra]
\label{prop:poly-relative-one-object-cascade-response}
In the one-object case, the relative Mealy response coalgebra is the total
response profile of the cascade action
\[
(x,y)\cdot m
=
\bigl(x\cdot m,\;y\cdot\omega_P(m,x)\bigr).
\]
\end{proposition}

\begin{proof}
\leavevmode
\begin{enumerate}
\item \textbf{Identify the state object.}
Since both object objects are terminal,
\[
S_{\Phi,Y}=P\times Y.
\]

\item \textbf{Identify the response interface.}
The interface is
\[
J_Y\cong Y.
\]

\item \textbf{Specialize the coalgebra.}
The general formula
\[
c_{\Phi,Y}(x,y)(a)
=
\left(
\omega_P(a,x),
\left(
\delta_P(a,x),
\sigma(\omega_P(a,x),y)
\right)
\right)
\]
becomes
\[
c_{\Phi,Y}(x,y)(m)
=
\left(
\omega_P(m,x),
\left(
x\cdot m,
y\cdot\omega_P(m,x)
\right)
\right).
\]
\end{enumerate}
\end{proof}

\subsection{One-object simulations}
\label{subsec:poly-relative-one-object-simulations}

A one-object Mealy simulation
\[
\Theta:\Phi\Rightarrow\Psi
\]
has components
\[
\theta:Q\to P,
\qquad
\alpha:Q\to N.
\]
For every right \(N\)-object \(Y\), the evaluated map is
\[
Q\times Y\to P\times Y,
\qquad
(y,z)\longmapsto(\theta y,z\cdot\alpha_y).
\]
The output-lax comparison component is
\[
\chi(y,z)=\alpha_y.
\]
The simulation equation becomes
\[
\alpha_y\,\omega_P(m,\theta y)
=
\omega_Q(m,y)\,\alpha_{y\cdot m}.
\]

\begin{proposition}[One-object output-lax profile comparison]
\label{prop:poly-relative-one-object-output-lax}
A one-object Mealy simulation evaluates to an output-lax profile comparison
between the corresponding cascade response coalgebras.
\end{proposition}

\begin{proof}
\leavevmode
\begin{enumerate}
\item \textbf{Identify the evaluated state map.}
The state map is
\[
(y,z)\mapsto(\theta y,z\cdot\alpha_y).
\]

\item \textbf{Identify strict input-target preservation.}
The state component of the simulation says
\[
\theta(y\cdot m)=\theta y\cdot m.
\]

\item \textbf{Identify output-lax naturality.}
The output equation
\[
\alpha_y\,\omega_P(m,\theta y)
=
\omega_Q(m,y)\,\alpha_{y\cdot m}
\]
is exactly the one-object output-lax naturality equation.

\item \textbf{Flatten.}
Flattening gives the input-strict functor between the one-object program
categories.
\end{enumerate}
\end{proof}

\section{The stateless specialization}
\label{sec:poly-relative-stateless-specialization}

\subsection{Stateless response coalgebras}
\label{subsec:poly-relative-stateless-response-coalgebras}

Let
\[
F:\mathbb A\to\mathbb B
\]
be an internal functor with object map \(F_0\) and arrow map \(F_1\).  Regard
\(F\) as the representable Mealy morphism with carrier
\[
\begin{tikzcd}
&
A_0
  \arrow[dl,"1_{A_0}"']
  \arrow[dr,"F_0"]
&
\\
A_0 && B_0 .
\end{tikzcd}
\]

For a right \(\mathbb B\)-action \(Y\), the relative state object is
\[
S_{F,Y}:=A_0\times_{F_0,B_0,r}Y.
\]
It is displayed by
\[
\begin{tikzcd}
S_{F,Y}
  \arrow[r,"\pi_Y"]
  \arrow[d,"\pi_A"']
&
Y
  \arrow[d,"r"]
\\
A_0
  \arrow[r,"F_0"']
&
B_0
\arrow[ul,phantom,very near start,"\lrcorner"] .
\end{tikzcd}
\]
The induced coalgebra is
\[
c_{F,Y}:
S_{F,Y}\to
\mathcal M_{\mathbb A,\mathbb B,Y}(S_{F,Y}),
\]
given by
\[
c_{F,Y}(v,y)(a)
=
\left(
F_1(a),
\left(
s_A(a),
\sigma(F_1(a),y)
\right)
\right),
\]
for every arrow \(a:u\to v\).

\begin{proposition}[Stateless relative response coalgebra]
\label{prop:poly-relative-stateless-response-coalgebra}
The coalgebra \(c_{F,Y}\) is the relative response coalgebra associated to the
representable Mealy morphism induced by \(F\).  Its flattened action category
is the action category of the ordinary restricted action \(F^\ast Y\).
\end{proposition}

\begin{proof}
\leavevmode
\begin{enumerate}
\item \textbf{Specialize the carrier.}
For the representable carrier, the state component is the current input object
\(v\in A_0\).

\item \textbf{Identify the state update.}
The state update is forced:
\[
\delta(a,v)=s_A(a).
\]

\item \textbf{Identify the output arrow.}
The emitted output arrow is
\[
\omega(a,v)=F_1(a).
\]

\item \textbf{Substitute into the general coalgebra formula.}
Substituting these into
Definition~\ref{def:poly-relative-mealy-response-coalgebra} gives the
displayed coalgebra.

\item \textbf{Flatten.}
Flattening recovers the usual restriction of a right \(\mathbb B\)-action
along \(F\).
\end{enumerate}
\end{proof}

\subsection{Terminal stateless case}
\label{subsec:poly-relative-terminal-stateless-case}

\begin{corollary}[Terminal stateless coalgebra]
\label{cor:poly-relative-terminal-stateless-coalgebra}
For
\[
Y=\mathbf 1_{\mathbb B},
\]
one has
\[
S_{F,\mathbf 1_{\mathbb B}}\cong A_0,
\]
and the coalgebra becomes
\[
c_F(v)(a)=(F_1(a),s_A(a)).
\]
The recovered program category is
\[
\mathbb A^{\mathrm{op}},
\]
and the output labelling is
\[
F^{\mathrm{op}}:\mathbb A^{\mathrm{op}}\to\mathbb B^{\mathrm{op}}.
\]
\end{corollary}

\begin{proof}
\leavevmode
\begin{enumerate}
\item \textbf{Identify the state object.}
For the terminal downstream action,
\[
S_{F,\mathbf 1_{\mathbb B}}
=
A_0\times_{F_0,B_0,1_{B_0}}B_0
\cong A_0.
\]

\item \textbf{Identify the coalgebra.}
The action of an input arrow \(a:u\to v\) sends \(v\) to \(u=s_A(a)\) and
emits \(F_1(a)\).  Hence
\[
c_F(v)(a)=(F_1(a),s_A(a)).
\]

\item \textbf{Flatten.}
The recovered program category is the action category of the terminal right
\(\mathbb A\)-action, namely
\[
\mathbb A^{\mathrm{op}}.
\]

\item \textbf{Identify the output labelling.}
The output labelling sends \(a\) to \(F_1(a)\), viewed oppositely, hence is
\[
F^{\mathrm{op}}.
\]
\end{enumerate}
\end{proof}

\section{Examples}
\label{sec:poly-relative-examples}

\subsection{Terminal downstream response polynomial}
\label{subsec:poly-relative-example-terminal-downstream-response}

Let
\[
Y=\mathbf 1_{\mathbb B}
=
(B_0\xrightarrow{1_{B_0}}B_0,\tau_B)
\]
be the terminal right \(\mathbb B\)-action.  Then
\[
\tau_B(\beta,t_B\beta)=s_B\beta.
\]
The response interface is
\[
J_{\mathbf 1_{\mathbb B}}=A_0\times B_0.
\]
The fibre formula becomes
\[
\mathcal M_{\mathbb A,\mathbb B,\mathbf 1_{\mathbb B}}(X)_{(v,b)}
\cong
\prod_{a:u\to v}
\sum_{\beta:b'\to b}
X_{(u,b')}.
\]

Moreover,
\[
S_{\Phi,\mathbf 1_{\mathbb B}}
=
P\times_{B_0}B_0
\cong P.
\]
Under this isomorphism,
\[
c_{\Phi,\mathbf 1_{\mathbb B}}(x)(a)
=
(\omega_P(a,x),\delta_P(a,x)).
\]

\begin{corollary}[Terminal Mealy polynomial]
\label{cor:poly-relative-terminal-mealy-polynomial}
The terminal downstream response polynomial of a Mealy morphism
\[
\Phi:\mathbb A\rightsquigarrow\mathbb B
\]
is the special case \(Y=\mathbf 1_{\mathbb B}\) of the relative response
polynomial.  Its induced coalgebra is the terminal output-response profile
\[
x
\longmapsto
\left(
a
\longmapsto(\omega_P(a,x),\delta_P(a,x))
\right).
\]
Flattening this coalgebra recovers the terminal relative program category
\[
\mathsf{Prog}_\Phi.
\]
\end{corollary}

\subsection{Terminal simulation comparison}
\label{subsec:poly-relative-example-terminal-simulation}

\begin{proposition}[Terminal simulation comparison]
\label{prop:poly-relative-terminal-simulation-comparison}
Let
\[
\Theta:\Phi\Rightarrow\Psi
\]
be a Mealy simulation with components \((\theta,\alpha)\).  In the terminal
downstream case, the output-lax profile comparison has state map
\[
\theta:Q\to P
\]
and output-comparison component
\[
\alpha_y:q_P(\theta y)\to q_Q(y).
\]
After flattening, it gives the input-strict functor
\[
\mathsf{Prog}_\Psi\to\mathsf{Prog}_\Phi
\]
and the natural transformation
\[
\Lambda_{\Psi,\mathbf 1_{\mathbb B}}
\Rightarrow
\Lambda_{\Phi,\mathbf 1_{\mathbb B}}H_\Theta
\]
with components
\[
\alpha_y^{\mathrm{op}}:
q_Q(y)\longrightarrow q_P(\theta y)
\]
in \(\mathbb B^{\mathrm{op}}\), where
\[
\alpha_y:q_P(\theta y)\to q_Q(y)
\]
is the underlying arrow of \(\mathbb B\).
\end{proposition}

\begin{proof}
\leavevmode
\begin{enumerate}
\item \textbf{Identify terminal evaluated states.}
In the terminal downstream action,
\[
S_{\Psi,\mathbf 1_{\mathbb B}}\cong Q
\qquad\text{and}\qquad
S_{\Phi,\mathbf 1_{\mathbb B}}\cong P.
\]
The evaluated state map is therefore
\[
\theta:Q\to P.
\]

\item \textbf{Identify the output comparison.}
The output-comparison component remains
\[
\alpha_y:q_P(\theta y)\to q_Q(y)
\]
in \(\mathbb B\).

\item \textbf{View the comparison oppositely.}
Since the terminal downstream action category is
\[
\int_{\mathbb B}\mathbf 1_{\mathbb B}\cong\mathbb B^{\mathrm{op}},
\]
the component is viewed as
\[
\alpha_y^{\mathrm{op}}:
q_Q(y)\to q_P(\theta y)
\]
in \(\mathbb B^{\mathrm{op}}\).

\item \textbf{Flatten.}
Flattening gives the input-strict functor between terminal relative program
categories and the displayed natural transformation.
\end{enumerate}
\end{proof}

\subsection{Word inputs and total response profiles}
\label{subsec:poly-relative-example-words}

If \(\mathbb A=I^\ast\) is the one-object category associated to the free
monoid of input words, then the response profile at a state is total over all
words:
\[
w\longmapsto
\left(
\omega_P(w,x),
\delta_P(w,x)
\right).
\]
Letter-step execution is obtained either by restricting one-step inputs to
the length-one subobject of \(I^\ast\), or by beginning with a generator-level
operation and extending it multiplicatively to \(I^\ast\).


\section{Chapter synthesis}
\label{sec:relative-response-synthesis}

\subsection{Operational evaluation}
\label{subsec:relative-response-operational-summary}

\begin{theorem}[Relative execution by downstream evaluation]
\label{thm:relative-response-operational-summary}
For every Mealy morphism
\[
    \Phi:\mathbb A\rightsquigarrow\mathbb B
\]
and right \(\mathbb B\)-action \(Y\), the following constructions agree:
\begin{enumerate}
    \item the action category of the restricted action \(\Phi^{\ast}Y\);
    \item the execution category of the composite
    \[
        \widehat Y\circ\Phi:\mathbb A\rightsquigarrow\mathbf 1;
    \]
    \item the labelled pullback
    \[
        \operatorname{Exec}(\Phi)
        \times_{\mathbb B^{\mathrm{op}}}
        \int_{\mathbb B}Y.
    \]
\end{enumerate}
The common internal category is \(\mathsf{Prog}_{\Phi,Y}\).  Its object
object is
\[
    S_{\Phi,Y}=P\times_{B_0}Y,
\]
and its one-step arrow object is
\[
    A_1\times_{A_0}S_{\Phi,Y}.
\]
\end{theorem}

\begin{proof}
The restricted action is Theorem
\ref{thm:relative-mealy-restriction-actions}.  Its action category is
Theorem~\ref{thm:relative-program-categories}.  The composite and pullback
presentations are Definition~\ref{def:relative-execution-category-composite}
and Theorem~\ref{thm:relative-execution-labelled-pullback}.
\end{proof}

\subsection{Currying and flattening}
\label{subsec:relative-response-currying-summary}

\begin{theorem}[Relative response equivalence]
\label{thm:relative-response-main-equivalence}
Assume the required dependent products exist.  Then:
\begin{enumerate}
    \item coalgebras
    \[
        c:X\to\mathcal M_{\mathbb A,\mathbb B,Y}(X)
    \]
    are equivalent to typed labelled one-step response maps;

    \item multiplicative coalgebras are equivalent to labelled relative
    actions
    \[
        \Lambda_X:\int_{\mathbb A}X\to\int_{\mathbb B}Y;
    \]

    \item the canonical Mealy coalgebra
    \[
        c_{\Phi,Y}(x,y)(a)
        =
        \left(
            \omega_P(a,x),
            \left(
                \delta_P(a,x),
                \sigma(\omega_P(a,x),y)
            \right)
        \right)
    \]
    is the dependent-product transpose of
    \((\Phi^{\ast}Y,\Lambda_{\Phi,Y})\);

    \item flattening \(c_{\Phi,Y}\) recovers
    \(\mathsf{Prog}_{\Phi,Y}\), including its input and downstream
    comparison functors;

    \item finite paths of the flattened category recover the relative
    trajectory pullback
    \[
        \mathscr T_\Phi
        \times_{\operatorname{Path}(\mathbb B^{\mathrm{op}})}
        \operatorname{Path}\!\left(\int_{\mathbb B}Y\right).
    \]
\end{enumerate}
\end{theorem}

\begin{proof}
The first two statements are Theorem
\ref{thm:poly-relative-currying-labelled-one-step} and Proposition
\ref{prop:poly-relative-coalgebras-labelled-actions}.  The induced coalgebra
and transpose theorem are Definition
\ref{def:poly-relative-mealy-response-coalgebra} and Theorem
\ref{thm:poly-relative-mealy-coalgebra-transpose}.  Flattening is Theorem
\ref{thm:poly-relative-program-category-recovered}.  The trajectory statement
is Corollary~\ref{cor:relative-mealy-response-trajectories}.
\end{proof}

\subsection{Variance and comparison structure}
\label{subsec:relative-response-variance-summary}

\begin{theorem}[Simulations after evaluation, currying, and flattening]
\label{thm:relative-response-simulation-summary}
Let
\[
    \Theta:\Phi\Rightarrow\Psi
\]
be a Mealy simulation.  For every right \(\mathbb B\)-action \(Y\):
\begin{enumerate}
    \item downstream evaluation gives an equivariant state map
    \[
        h_{\Theta,Y}:S_{\Psi,Y}\to S_{\Phi,Y};
    \]

    \item this state map induces a contravariant internal functor
    \[
        H_{\Theta,Y}:
        \mathsf{Prog}_{\Psi,Y}	o\mathsf{Prog}_{\Phi,Y};
    \]

    \item the raw output comparison gives a natural transformation
    \[
        \Lambda_{\Psi,Y}
        \Rightarrow
        \Lambda_{\Phi,Y}H_{\Theta,Y};
    \]

    \item polynomially, the same data form an output-lax profile comparison
    \[
        c_{\Psi,Y}\longrightarrow c_{\Phi,Y};
    \]

    \item flattening this output-lax comparison recovers the functor and
    natural transformation above.
\end{enumerate}
\end{theorem}

\begin{proof}
Operational evaluation is Theorem
\ref{thm:relative-simulations-evaluate-equivariantly}; the induced program
functor and downstream transformation are Theorem
\ref{thm:relative-simulations-program-functors} and Theorem
\ref{thm:relative-output-comparison-natural}.  The polynomial comparison and
its flattening are Theorem
\ref{thm:poly-relative-simulations-profile-comparisons}.
\end{proof}

\begin{remark}[Conceptual summary]
\label{rem:relative-response-final-conceptual-summary}
No new machine is introduced by the polynomial construction.  The labelled
relative action, multiplicative response coalgebra, flattened program
category, and finite relative trajectory object are four presentations of the
same evaluated Mealy semantics:
\[
\begin{tikzcd}[column sep=large,row sep=large]
    (\Phi^{\ast}Y,\Lambda_{\Phi,Y})
        \arrow[r,shift left=1.2ex,"\text{curry}"]
        \arrow[d,"\int_{\mathbb A}"']
    &
    c_{\Phi,Y}
        \arrow[l,shift left=1.2ex,"\text{uncurry}"']
        \arrow[d,"\text{flatten}"]
    \\
    \mathsf{Prog}_{\Phi,Y}
        \arrow[r,equal]
        \arrow[d,"\operatorname{Path}"']
    &
    \mathsf{Prog}_{\Phi,Y}
        \arrow[d,"\operatorname{Path}"]
    \\
    \mathscr T_{\Phi,Y}
        \arrow[r,equal]
    &
    \mathscr T_{\Phi,Y}.
\end{tikzcd}
\]
The next chapter places predicates on each of these presentations.  Span
modalities quantify over selected one-step execution witnesses, polynomial
predicate liftings quantify over positions in total response profiles, and
trajectory modalities quantify over the finite paths of the relative
execution category.
\end{remark}
